\section{El teorema fundamental de los morfismos propios. Aplicaciones}
\label{section:III.4}


\subsection{El teorema fundamental}
\label{subsection:III.4.1}

\begin{env}[4.1.1]
\label{III.4.1.1}
Sean $X$ e $Y$ esquemas ordinarios noetherianos, sea $f:X\to Y$ un
morfismo propio, sea $Y'$ un subconjunto cerrado de $Y$, y pongamos
\[
  X'=f^{-1}(Y').
\]
Denotamos por $\widehat{X}$ y $\widehat{Y}$ los esquemas formales
$X_{/X'}$ y $Y_{/Y'}$, obtenidos completando $X$ e $Y$ a lo largo de
$X'$ e $Y'$, respectivamente \sref[I]{I.10.8.5}.
Denotamos por $\widehat{f}:\widehat{X}\to\widehat{Y}$ la extensión de
$f$ a estas completaciones \sref[I]{I.10.9.1}.
Para todo $\sh{O}_X$-módulo coherente $\sh{F}$, denotamos por

\oldpage[III]{123}

\[
  \widehat{\sh{F}}=\sh{F}_{/X'}
\]
su completación a lo largo de $X'$ \sref[I]{I.10.8.4}; este es un
$\sh{O}_{\widehat{X}}$-módulo coherente \sref[I]{I.10.8.8}.
\end{env}

\begin{env}[4.1.2]
\label{III.4.1.2}
Sea $\sh{J}$ un ideal coherente de $\sh{O}_Y$ tal que
\[
  \Supp(\sh{O}_Y/\sh{J})=Y'
\]
\sref[I]{I.5.2.1}.
Sabemos \sref[I]{I.4.4.5} que
\[
  \sh{K}=f^*(\sh{J})\sh{O}_X
\]
es un ideal coherente de $\sh{O}_X$ que satisface
\[
  \Supp(\sh{O}_X/\sh{K})=X'.
\]
Para todo $k\geq0$, consideremos el $\sh{O}_X$-módulo coherente
\[
  \sh{F}_k
  =
  \sh{F}\otimes_{\sh{O}_Y}(\sh{O}_Y/\sh{J}^{k+1})
  =
  \sh{F}/\sh{K}^{k+1}\sh{F}.
\]
Los $\sh{O}_Y$-módulos
\[
  \RR^n f_*(\sh{F})
  \quad\text{y}\quad
  \RR^n f_*(\sh{F}_k)
\]
son coherentes para todo $n$ \sref{III.3.2.1}.
Para todo $k\geq0$ y todo $n$, el homomorfismo canónico
$\sh{F}\to\sh{F}_k$ induce, por funtorialidad, un homomorfismo
\[
  \RR^n f_*(\sh{F})
  \longrightarrow
  \RR^n f_*(\sh{F}_k).
  \tag{4.1.2.1}
  \label{III.4.1.2.1}
\]
Como $\sh{F}_k$ es un $\sh{O}_X/\sh{K}^{k+1}$-módulo,
$\RR^n f_*(\sh{F}_k)$ es un
$\sh{O}_Y/\sh{J}^{k+1}$-módulo \sref[0]{0.12.2.1}.
Deducimos, pues, de \eqref{III.4.1.2.1} un homomorfismo
\[
  \RR^n f_*(\sh{F})
  \otimes_{\sh{O}_Y}(\sh{O}_Y/\sh{J}^{k+1})
  \longrightarrow
  \RR^n f_*(\sh{F}_k).
  \tag{4.1.2.2}
  \label{III.4.1.2.2}
\]
Los dos miembros de \eqref{III.4.1.2.2} forman sistemas proyectivos.
El límite proyectivo del miembro izquierdo es la completación
\[
  \bigl(\RR^n f_*(\sh{F})\bigr)_{/Y'},
\]
que denotaremos por
\[
  \widehat{\RR^n f_*(\sh{F})}.
\]
Los homomorfismos \eqref{III.4.1.2.2} son compatibles con los sistemas
proyectivos y definen, por tanto, al pasar al límite, un homomorfismo
canónico
\[
  \vphi_n:
  \widehat{\RR^n f_*(\sh{F})}
  \longrightarrow
  \varprojlim_k\RR^n f_*(\sh{F}_k).
  \tag{4.1.2.3}
  \label{III.4.1.2.3}
\]
Además, \eqref{III.4.1.2.2} es un homomorfismo de
$\sh{O}_Y/\sh{J}^{k+1}$-módulos.
Por \sref[I]{I.10.8.3}, puede considerarse, por tanto, como un
homomorfismo continuo de $\sh{O}_{\widehat{Y}}$-módulos topológicos
provistos de sus topologías seudodiscretas \sref[0]{0.3.8.1}.
En consecuencia, $\vphi_n$ es un homomorfismo continuo de
$\sh{O}_{\widehat{Y}}$-módulos topológicos.
\end{env}

\begin{env}[4.1.3]
\label{III.4.1.3}
Sea $i:\widehat{X}\to X$ el morfismo canónico de espacios anillados
definido en \sref[I]{I.10.8.7}.
Tenemos el diagrama conmutativo
\[
  \xymatrix{
    X_k\ar[r]^{h_k}\ar[dr]_{i_k} &
    \widehat{X}\ar[d]^{i}\\
    & X
  }
  \tag{4.1.3.1}
  \label{III.4.1.3.1}
\]
donde $X_k$ es el subesquema cerrado de $X$ definido por
$\sh{K}^{k+1}$, $i_k$ es la inmersión canónica, y $h_k$ es el
morfismo de espacios anillados que corresponde a la identidad sobre
los espacios subyacentes y al homomorfismo canónico
\[
  \sh{O}_{\widehat{X}}
  \longrightarrow
  \sh{O}_X/\sh{K}^{k+1}
\]
\sref[I]{I.10.5.2}.
Además,
\[
  \widehat{\sh{F}}=i^*(\sh{F})
\]
salvo isomorfismo canónico \sref[I]{I.10.8.8}.
Tenemos
\[
  \HH^n(X_k,i_k^*(\sh{F}_k))
  =
  \HH^n(X,\sh{F}_k)
  \tag{4.1.3.2}
  \label{III.4.1.3.2}
\]
salvo isomorfismo canónico.

\oldpage[III]{124}

En efecto,
\[
  \sh{F}_k=(i_k)_*(i_k^*(\sh{F}_k))
\]
(G,~II,~4.9.1).
El homomorfismo canónico
\[
  \HH^n(\widehat{X},\widehat{\sh{F}})
  \longrightarrow
  \HH^n(X_k,h_k^*(\widehat{\sh{F}}))
\]
\sref[0]{0.12.1.3.5} puede escribirse también
\[
  \HH^n(\widehat{X},\widehat{\sh{F}})
  \longrightarrow
  \HH^n(X,\sh{F}_k).
  \tag{4.1.3.3}
  \label{III.4.1.3.3}
\]
Estos homomorfismos forman un sistema proyectivo y definen, por tanto,
un homomorfismo canónico
\[
  \psi_{n,X}:
  \HH^n(\widehat{X},\widehat{\sh{F}})
  \longrightarrow
  \varprojlim_k\HH^n(X,\sh{F}_k).
  \tag{4.1.3.4}
  \label{III.4.1.3.4}
\]
Reemplazando $X$ por un subconjunto abierto de la forma $f^{-1}(V)$,
donde $V$ es un abierto afín de $Y$, y teniendo en cuenta
\sref{III.1.4.11}, obtenemos homomorfismos
\[
  \psi_{n,V}:
  \HH^n(\widehat{X}\cap f^{-1}(V),\widehat{\sh{F}})
  \longrightarrow
  \varprojlim_k\Gamma(V,\RR^n f_*(\sh{F}_k)).
  \tag{4.1.3.5}
  \label{III.4.1.3.5}
\]
Estos homomorfismos conmutan manifiestamente con la restricción de $V$
a un subconjunto abierto afín más pequeño.
Definen, por tanto, un homomorfismo canónico de haces
\[
  \psi_n:
  \RR^n\widehat{f}_*(\widehat{\sh{F}})
  \longrightarrow
  \varprojlim_k\RR^n f_*(\sh{F}_k).
  \tag{4.1.3.6}
  \label{III.4.1.3.6}
\]
\end{env}

\begin{env}[4.1.4]
\label{III.4.1.4}
Por último, sea $j:\widehat{Y}\to Y$ el morfismo canónico de espacios
anillados \sref[I]{I.10.8.7}.
Como $\RR^n f_*(\sh{F})$ es un $\sh{O}_Y$-módulo coherente
\sref{III.3.2.1}, tenemos, salvo isomorfismo canónico,
\[
  j^*(\RR^n f_*(\sh{F}))
  =
  \widehat{\RR^n f_*(\sh{F})}
\]
\sref[I]{I.10.8.8}.
Existe, en consecuencia, un homomorfismo canónico
\[
  \rho_n:
  \widehat{\RR^n f_*(\sh{F})}
  =
  j^*(\RR^n f_*(\sh{F}))
  \longrightarrow
  \RR^n\widehat{f}_*(i^*(\sh{F}))
  =
  \RR^n\widehat{f}_*(\widehat{\sh{F}})
  \tag{4.1.4.1}
  \label{III.4.1.4.1}
\]
definido para los espacios anillados en general; véase la demostración
de \sref{III.1.4.15}.
Demostraremos que el diagrama
\[
  \xymatrix{
    \widehat{\RR^n f_*(\sh{F})}
      \ar[rr]^{\rho_n}\ar[dr]_{\vphi_n}
    &&
    \RR^n\widehat{f}_*(\widehat{\sh{F}})
      \ar[dl]^{\psi_n}\\
    &
    \displaystyle\varprojlim_k\RR^n f_*(\sh{F}_k)
    &
  }
  \tag{4.1.4.2}
  \label{III.4.1.4.2}
\]
es conmutativo.
Basta demostrar la conmutatividad del diagrama formado por los
homomorfismos correspondientes de prehaces.
Podemos, por tanto, restringirnos al caso en que $Y$ es afín, y basta
demostrar que
\[
  \xymatrix{
    \widehat{\HH^n(X,\sh{F})}
      \ar[rr]^{\rho_n}\ar[dr]_{\vphi_n}
    &&
    \HH^n(\widehat{X},\widehat{\sh{F}})
      \ar[dl]^{\psi_{n,X}}\\
    &
    \displaystyle\varprojlim_k\HH^n(X,\sh{F}_k)
    &
  }
  \tag{4.1.4.3}
  \label{III.4.1.4.3}
\]
es conmutativo.
La conmutatividad de \eqref{III.4.1.3.1}, junto con las relaciones
entre los grupos de cohomología establecidas en \sref{III.4.1.3}, da
el diagrama conmutativo

\oldpage[III]{125}

\[
  \xymatrix{
    \HH^n(X,\sh{F})
      \ar[r]\ar[dr]
    &
    \HH^n(\widehat{X},\widehat{\sh{F}})
      =
    \HH^n(\widehat{X},i^*(\sh{F}))
      \ar[d]\\
    &
    \HH^n(X,\sh{F}_k)
      =
    \HH^n(X_k,i_k^*(\sh{F})).
  }
\]
La conmutatividad de \eqref{III.4.1.4.3} se sigue inmediatamente.
\end{env}

\begin{theorem}[4.1.5]
\label{III.4.1.5}
Sea $f:X\to Y$ un morfismo propio de esquemas noetherianos, sea
$Y'$ un subconjunto cerrado de $Y$, y pongamos $X'=f^{-1}(Y')$.
Entonces, para todo $\sh{O}_X$-módulo coherente $\sh{F}$,
\[
  \RR^n\widehat{f}_*(\widehat{\sh{F}})
\]
es un $\sh{O}_{\widehat{X}}$-módulo coherente, y los homomorfismos
$\vphi_n$, $\psi_n$ y $\rho_n$ del diagrama
\eqref{III.4.1.4.2} son isomorfismos topológicos.
\end{theorem}

\begin{proof}
Basta demostrar que $\vphi_n$ y $\psi_n$ son isomorfismos.
En efecto, $\RR^n f_*(\sh{F})$ es coherente \sref{III.3.2.1}, luego
su completación $\widehat{\RR^n f_*(\sh{F})}$ es coherente
\sref[I]{I.10.8.8}; la bicontinuidad de $\vphi_n$, $\psi_n$ y
$\rho_n$ es entonces automática \sref[I]{I.10.11.6}.
\end{proof}

\begin{env}[4.1.6]
\label{III.4.1.6}
\emph{Observaciones.}
\begin{enumerate}
  \item[(i)] Pongamos
  \[
    \widehat{\sh{F}}_k
    =
    \widehat{\sh{F}}/
    \widehat{\sh{K}}^{\,k+1}\widehat{\sh{F}}.
  \]
  Entonces, inmediatamente,
  \[
    \sh{F}_k=i_*(\widehat{\sh{F}}_k),
  \]
  y el homomorfismo canónico \eqref{III.4.1.3.6} es precisamente el
  homomorfismo ya definido en (\hyperref[III.3.4.2.2]{3.4.2.2}):
  \[
    \RR^n\widehat{f}_*(\widehat{\sh{F}})
    \longrightarrow
    \varprojlim_k
    \RR^n\widehat{f}_*(\widehat{\sh{F}}_k).
    \tag{4.1.6.1}
    \label{III.4.1.6.1}
  \]
  Así, el hecho de que $\psi_n$ sea un isomorfismo es un caso particular
  de \sref{III.3.4.3}.
  Daremos, no obstante, más adelante una demostración directa, evitando
  las delicadas consideraciones relativas a los límites proyectivos de
  sucesiones espectrales \sref[0]{0.13.7} en las que se apoya el teorema
  general \sref{III.3.4.3}.

  \item[(ii)] Como los $\psi_n$ son isomorfismos, decir que los
  $\vphi_n$ son isomorfismos equivale a decir que
  \[
    \rho_n=\psi_n^{-1}\circ\vphi_n
  \]
  son isomorfismos.
  El Teorema~\sref{III.4.1.5} afirma, en particular, que la formación de
  $\RR^n f_*$ conmuta con la completación; puede llamarse, por tanto, el
  primer teorema de comparación entre las teorías ``algebraica'' y
  ``formal''.
\end{enumerate}
Comenzaremos demostrando la forma afín de \sref{III.4.1.5}.
\end{env}

\begin{corollary}[4.1.7]
\label{III.4.1.7}
Bajo las hipótesis de \sref{III.4.1.5}, supongamos además que
$Y=\Spec(A)$, donde $A$ es noetheriano, y que
$\sh{J}=\widetilde{\mathfrak{J}}$, donde $\mathfrak{J}$ es un ideal
de $A$, de modo que
\[
  \sh{F}_k=\sh{F}/\mathfrak{J}^{k+1}\sh{F}.
\]
El homomorfismo canónico
\[
  \vphi_n:
  \widehat{\HH^n(X,\sh{F})}
  \longrightarrow
  \varprojlim_k\HH^n(X,\sh{F}_k)
  \tag{4.1.7.1}
  \label{III.4.1.7.1}
\]
donde el primer término es la completación separada de
$\HH^n(X,\sh{F})$ para la topología $\mathfrak{J}$-preádica, es un
isomorfismo.
Para todo $n$, el sistema proyectivo
\[
  \bigl(\HH^n(X,\sh{F}_k)\bigr)_{k\geq0}
\]
satisface la condición \emph{(ML)}, y el homomorfismo canónico
\[
  \psi_n:
  \HH^n(\widehat{X},\widehat{\sh{F}})
  \longrightarrow
  \varprojlim_k\HH^n(X,\sh{F}_k)
  \tag{4.1.7.2}
  \label{III.4.1.7.2}
\]

\oldpage[III]{126}

es un isomorfismo.
Por último, la filtración sobre $\HH^n(X,\sh{F})$ definida por los
núcleos de los homomorfismos canónicos
\[
  \HH^n(X,\sh{F})\longrightarrow\HH^n(X,\sh{F}_k)
\]
es $\mathfrak{J}$-buena \sref[0]{0.13.7.7}, y $\vphi_n$ es un
isomorfismo topológico.\footnote{La demostración siguiente, más sencilla
que la demostración original, y el complemento relativo a la filtración
de $\HH^n(X,\sh{F})$, nos fueron comunicados por J.-P. Serre.}
\end{corollary}

\begin{proof}
Fijemos $n\geq0$ durante toda esta demostración y, para simplificar,
pongamos
\[
  \mathrm{H}=\HH^n(X,\sh{F}),
  \qquad
  \mathrm{H}_k=\HH^n(X,\sh{F}_k).
  \tag{4.1.7.3}
  \label{III.4.1.7.3}
\]
Pongamos también
\[
  \mathrm{R}_k
  =
  \Ker(\mathrm{H}\longrightarrow\mathrm{H}_k),
  \qquad
  \text{un sub-$A$-módulo de $\mathrm{H}$.}
  \tag{4.1.7.4}
  \label{III.4.1.7.4}
\]
La sucesión exacta de cohomología
\[
  \HH^n(X,\mathfrak{J}^{k+1}\sh{F})
  \longrightarrow
  \HH^n(X,\sh{F})
  \longrightarrow
  \HH^n(X,\sh{F}_k)
  \xrightarrow{\partial}
  \HH^{n+1}(X,\mathfrak{J}^{k+1}\sh{F})
  \longrightarrow
  \HH^{n+1}(X,\sh{F})
\]
muestra que
\[
  \mathrm{R}_k
  =
  \Im\bigl(
    \HH^n(X,\mathfrak{J}^{k+1}\sh{F})
    \longrightarrow
    \HH^n(X,\sh{F})
  \bigr).
\]
Ponemos
\[
  \begin{split}
  \mathrm{Q}_k
  &=
  \Ker\bigl(
    \HH^{n+1}(X,\mathfrak{J}^{k+1}\sh{F})
    \longrightarrow
    \HH^{n+1}(X,\sh{F})
  \bigr)
  \\
  &=
  \Im\bigl(
    \HH^n(X,\sh{F}_k)
    \longrightarrow
    \HH^{n+1}(X,\mathfrak{J}^{k+1}\sh{F})
  \bigr).
  \end{split}
  \tag{4.1.7.5}
  \label{III.4.1.7.5}
\]
Tenemos, por tanto, la sucesión exacta
\[
  0\longrightarrow\mathrm{R}_k
  \longrightarrow\mathrm{H}
  \longrightarrow\mathrm{H}_k
  \longrightarrow\mathrm{Q}_k
  \longrightarrow0.
  \tag{4.1.7.6}
  \label{III.4.1.7.6}
\]

\begin{env}[4.1.7.7]
\label{III.4.1.7.7}
Sea $x\in\mathfrak{J}^m$, donde $m\geq0$.
La multiplicación por $x$ sobre $\mathfrak{J}^k\sh{F}$ es un
homomorfismo
\[
  \mathfrak{J}^k\sh{F}
  \longrightarrow
  \mathfrak{J}^{k+m}\sh{F},
\]
y da, por tanto, un homomorfismo
\[
  \mu_{x,m}:
  \HH^n(X,\mathfrak{J}^k\sh{F})
  \longrightarrow
  \HH^n(X,\mathfrak{J}^{k+m}\sh{F}).
  \tag{4.1.7.8}
  \label{III.4.1.7.8}
\]
Sea
\[
  S=\bigoplus_{k\geq0}\mathfrak{J}^k
\]
la $A$-álgebra graduada.
Sabemos que las multiplicaciones $\mu_{x,m}$ dotan a
\[
  \mathrm{E}
  =
  \bigoplus_{k\geq0}
  \HH^n(X,\mathfrak{J}^k\sh{F})
\]
de una estructura de $S$-módulo graduado de tipo finito
\sref{III.3.3.2}; el anillo graduado $S$ es noetheriano
\sref[II]{II.2.1.5}.
\end{env}

\begin{lemma}[4.1.7.9]
\label{III.4.1.7.9}
Los submódulos $\mathrm{R}_k$ de $\mathrm{H}$ definen una
filtración $\mathfrak{J}$-buena sobre $\mathrm{H}$.
\end{lemma}

\begin{proof}
Demostremos primero que
\[
  \mathfrak{J}^m\mathrm{R}_k
  \subset
  \mathrm{R}_{k+m},
  \tag{4.1.7.10}
  \label{III.4.1.7.10}
\]
donde la multiplicación sobre
$\mathrm{H}=\HH^n(X,\sh{F})$ por un elemento
$x\in\mathfrak{J}^m$ es la aplicación $\mu_{x,0}$.
Para todo $x\in\mathfrak{J}^m$, el diagrama
\[
  \xymatrix{
    \mathfrak{J}^{k+1}\sh{F}
      \ar[r]^{x}\ar[d]
    &
    \mathfrak{J}^{k+m+1}\sh{F}
      \ar[d]
    \\
    \sh{F}
      \ar[r]_{x}
    &
    \sh{F}
  }
  \tag{4.1.7.11}
  \label{III.4.1.7.11}
\]

\oldpage[III]{127}

es conmutativo; sus flechas horizontales son las multiplicaciones por
$x$, y sus flechas verticales son las inyecciones canónicas.
En consecuencia, el diagrama correspondiente
\[
  \xymatrix{
    \HH^n(X,\mathfrak{J}^{k+1}\sh{F})
      \ar[r]^{\mu_{x,m}}\ar[d]
    &
    \HH^n(X,\mathfrak{J}^{k+m+1}\sh{F})
      \ar[d]
    \\
    \HH^n(X,\sh{F})
      \ar[r]_{\mu_{x,0}}
    &
    \HH^n(X,\sh{F})
  }
\]
es conmutativo.
La interpretación de $\mathrm{R}_k$ como la imagen de
\[
  \HH^n(X,\mathfrak{J}^{k+1}\sh{F})
  \longrightarrow
  \HH^n(X,\sh{F})
\]
demuestra, por tanto, \eqref{III.4.1.7.10}.
Muestra también que el $S$-módulo graduado
\[
  \mathrm{R}=\bigoplus_{k\geq0}\mathrm{R}_k
\]
es un cociente del sub-$S$-módulo
\[
  \mathrm{M}
  =
  \bigoplus_{k\geq0}
  \HH^n(X,\mathfrak{J}^{k+1}\sh{F})
\]
de $\mathrm{E}$.
La observación anterior muestra, pues, que $\mathrm{R}$ es un
$S$-módulo de tipo finito.
Esto equivale a la afirmación del lema
(Bourbaki, \emph{Alg.\ comm.}, cap.~III, \S3, n.º~1, teorema~1).
\end{proof}

\begin{env}[4.1.7.12]
\label{III.4.1.7.12}
Consideremos ahora el $S$-módulo graduado
\[
  \mathrm{N}
  =
  \bigoplus_{k\geq0}
  \HH^{n+1}(X,\mathfrak{J}^{k+1}\sh{F}),
\]
definido como en \eqref{III.4.1.7.8}.
También es un $S$-módulo de tipo finito por \sref{III.3.3.2}.
Por \eqref{III.4.1.7.5}, tenemos
$\mathrm{Q}_k\subset\mathrm{N}_k$ para todo $k$, y el diagrama
\eqref{III.4.1.7.11}, sustituyendo $n$ por $n+1$, muestra que
\[
  S_m\mathrm{Q}_k
  =
  \mathfrak{J}^m\mathrm{Q}_k
  \subset
  \mathrm{Q}_{k+m}.
\]
En otros términos, $\mathrm{Q}$ es un sub-$S$-módulo graduado de
$\mathrm{N}$ y es, por tanto, de tipo finito.
\end{env}

\begin{env}[4.1.7.13]
\label{III.4.1.7.13}
Denotemos por $\alpha_m$ la inyección canónica
\[
  \mathfrak{J}^m\longrightarrow A,
\]
que también puede escribirse $S_m\to S_0$.
Como $\mathfrak{J}^{k+1}\sh{F}_k=0$, el $A$-módulo
$\HH^n(X,\sh{F}_k)$ es anulado por $\mathfrak{J}^{k+1}$.
Como $\mathrm{Q}_k$ es la imagen del $A$-homomorfismo
\[
  \HH^n(X,\sh{F}_k)
  \longrightarrow
  \HH^{n+1}(X,\mathfrak{J}^{k+1}\sh{F}),
\]
$\mathrm{Q}_k$ es igualmente anulado por $\mathfrak{J}^{k+1}$.
Equivalentemente, en el $S$-módulo $\mathrm{Q}$,
\[
  \alpha_{k+1}(S_{k+1})\mathrm{Q}_k=0.
  \tag{4.1.7.14}
  \label{III.4.1.7.14}
\]
Como $\mathrm{Q}$ es un $S$-módulo de tipo finito, existen enteros
$k_0$ y $h$ tales que
\[
  \mathrm{Q}_{k+h}=S_h\mathrm{Q}_k
  \qquad(k\geq k_0)
\]
\sref[II]{II.2.1.6}[ii].
Esta relación y \eqref{III.4.1.7.14} implican que existe un entero
$r>0$ tal que
\[
  \alpha_r(S_r)\mathrm{Q}=0.
  \tag{4.1.7.15}
  \label{III.4.1.7.15}
\]
\end{env}

\begin{env}[4.1.7.16]
\label{III.4.1.7.16}
La inyección canónica
\[
  \mathfrak{J}^{k+m}\sh{F}
  \longrightarrow
  \mathfrak{J}^k\sh{F}
\]
induce en cohomología un $A$-homomorfismo
\[
  \nu_m:
  \HH^{n+1}(X,\mathfrak{J}^{k+m}\sh{F})
  \longrightarrow
  \HH^{n+1}(X,\mathfrak{J}^k\sh{F}).
  \tag{4.1.7.17}
  \label{III.4.1.7.17}
\]
Para todo $x\in\mathfrak{J}^m$, tenemos manifiestamente la
factorización
\[
  \xymatrix@C=4.5em{
    \HH^{n+1}(X,\mathfrak{J}^k\sh{F})
      \ar[r]^{\mu_{x,m}}
    &
    \HH^{n+1}(X,\mathfrak{J}^{k+m}\sh{F})
      \ar[r]^{\nu_m}
    &
    \HH^{n+1}(X,\mathfrak{J}^k\sh{F}),
  }
  \tag{4.1.7.18}
  \label{III.4.1.7.18}
\]
cuya compuesta es $\mu_{x,0}$.
\end{env}

\oldpage[III]{128}

Se sigue que, para todo sub-$A$-módulo $\mathrm{P}$ de
$\HH^{n+1}(X,\mathfrak{J}^k\sh{F})$, tenemos en el $S$-módulo
$\mathrm{N}$
\[
  \nu_m(S_m\mathrm{P})
  =
  \alpha_m(S_m)\mathrm{P}.
  \tag{4.1.7.19}
  \label{III.4.1.7.19}
\]

\begin{lemma}[4.1.7.20]
\label{III.4.1.7.20}
Existe un entero $m>0$ tal que
\[
  \nu_m(\mathrm{Q}_{k+m})=0
\]
para todo $k\geq k_0$.
\end{lemma}

\begin{proof}
Tomemos $m$ múltiplo de $h$ con $m\geq r$.
Como
\[
  \mathrm{Q}_{k+m}=S_m\mathrm{Q}_k
  \qquad(k\geq k_0),
\]
las relaciones \eqref{III.4.1.7.19} y \eqref{III.4.1.7.15} dan
\[
  \nu_m(\mathrm{Q}_{k+m})
  =
  \alpha_m(S_m)\mathrm{Q}_k
  \subset
  \alpha_r(S_r)\mathrm{Q}_k
  =
  0.
\]
\end{proof}

\begin{env}[4.1.7.21]
\label{III.4.1.7.21}
El diagrama conmutativo
\[
  \xymatrix@C=2.5em{
    \HH^n(X,\sh{F})
      \ar[r]
    &
    \HH^n(X,\sh{F}_k)
      \ar[r]^-{\partial}
    &
    \HH^{n+1}(X,\mathfrak{J}^{k+1}\sh{F})
      \ar[r]
    &
    \HH^{n+1}(X,\sh{F})
    \\
    \HH^n(X,\sh{F})
      \ar[r]\ar[u]
    &
    \HH^n(X,\sh{F}_{k+m})
      \ar[r]^-{\partial}\ar[u]
    &
    \HH^{n+1}(X,\mathfrak{J}^{k+m+1}\sh{F})
      \ar[r]\ar[u]
    &
    \HH^{n+1}(X,\sh{F})
      \ar[u]
  }
\]
procede a su vez del diagrama conmutativo
\[
  \xymatrix{
    0\ar[r]
    &
    \mathfrak{J}^{k+1}\sh{F}\ar[r]
    &
    \sh{F}\ar[r]
    &
    \sh{F}_k\ar[r]
    &
    0
    \\
    0\ar[r]
    &
    \mathfrak{J}^{k+m+1}\sh{F}\ar[r]\ar[u]
    &
    \sh{F}\ar[r]\ar[u]
    &
    \sh{F}_{k+m}\ar[r]\ar[u]
    &
    0.
  }
\]
Aquí las flechas verticales son las aplicaciones canónicas.
Obtenemos, por tanto, el diagrama conmutativo
\[
  \xymatrix@C=2.2em{
    0\ar[r]
    &
    \mathrm{R}_k\ar[r]
    &
    \mathrm{H}\ar[r]
    &
    \mathrm{H}_k\ar[r]
    &
    \mathrm{Q}_k\ar[r]
    &
    0
    \\
    0\ar[r]
    &
    \mathrm{R}_{k+m}\ar[r]\ar[u]
    &
    \mathrm{H}\ar[r]\ar@{=}[u]
    &
    \mathrm{H}_{k+m}\ar[r]\ar[u]
    &
    \mathrm{Q}_{k+m}\ar[r]\ar[u]^{\nu_m}
    &
    0,
  }
\]
de filas exactas.
Para $k\geq k_0$, la última flecha vertical es nula por
\sref{III.4.1.7.20}.
Así, la imagen de $\mathrm{H}_{k+m}$ en $\mathrm{H}_k$ está contenida
en
\[
  \Ker(\mathrm{H}_k\longrightarrow\mathrm{Q}_k)
  =
  \Im(\mathrm{H}\longrightarrow\mathrm{H}_k).
\]
Recíprocamente, la conmutatividad muestra que contiene
$\Im(\mathrm{H}\to\mathrm{H}_k)$ y, por tanto, que es igual a esta
imagen.
Lo mismo ocurre, pues, con las imágenes en $\mathrm{H}_k$ de
$\mathrm{H}_{k'}$ para $k'\geq k+m$.
Esto demuestra la condición \emph{(ML)} para el sistema proyectivo
$(\mathrm{H}_k)_{k\geq0}$.

Además, para todo subconjunto abierto afín $U$ de $X$,
\[
  \HH^i(U,\sh{F}_k)=0
  \qquad(i>0)
\]
por \sref{III.1.3.1}, y, para $m>0$, la aplicación
\[
  \HH^0(U,\sh{F}_{k+m})
  \longrightarrow
  \HH^0(U,\sh{F}_k)
\]
es sobreyectiva \sref[I]{I.1.3.9}.
Podemos, por tanto,

\oldpage[III]{129}

aplicar \sref[0]{0.13.3.1}; el homomorfismo canónico
\[
  \HH^n(\widehat{X},\widehat{\sh{F}})
  \longrightarrow
  \varprojlim_k\HH^n(X,\sh{F}_k)
\]
es biyectivo para todo $n\geq0$.

Como el sistema proyectivo
$(\mathrm{H}/\mathrm{R}_k)_{k\geq0}$ es estricto, podemos pasar al
límite proyectivo en las sucesiones exactas
\[
  0
  \longrightarrow
  \mathrm{H}/\mathrm{R}_k
  \longrightarrow
  \mathrm{H}_k
  \longrightarrow
  \mathrm{Q}_k
  \longrightarrow
  0
  \tag{4.1.7.22}
  \label{III.4.1.7.22}
\]
por \sref[0]{0.13.2.2}.
Como $\nu_m(\mathrm{Q}_{k+m})=0$, tenemos
\[
  \varprojlim_k\mathrm{Q}_k=0,
\]
y, por tanto, un isomorfismo topológico
\[
  \varprojlim_k(\mathrm{H}/\mathrm{R}_k)
  \isoto
  \varprojlim_k\mathrm{H}_k.
\]
Pero la filtración $(\mathrm{R}_k)$ de $\mathrm{H}$ es
$\mathfrak{J}$-buena y define, por tanto, sobre $\mathrm{H}$ la
topología $\mathfrak{J}$-preádica.
Así, $\varprojlim_k(\mathrm{H}/\mathrm{R}_k)$ es la completación
separada de $\mathrm{H}$ para esta topología.
Esto termina la demostración de \sref{III.4.1.7}.
\end{env}
\end{proof}

\begin{env}[4.1.8]
\label{III.4.1.8}
Podemos terminar ahora la demostración de \sref{III.4.1.5}.
Para todo subconjunto abierto afín $V$ de $Y$,
\[
  \Gamma\bigl(
    V,\widehat{\RR^n f_*(\sh{F})}
  \bigr)
\]
es la completación separada de
$\Gamma(V,\RR^n f_*(\sh{F}))$ para la topología
$\mathfrak{J}$-preádica, si
$\sh{J}|V=\widetilde{\mathfrak{J}}$, porque
$\RR^n f_*(\sh{F})$ es un $\sh{O}_Y$-módulo coherente
\sref[I]{I.10.8.4}.
Además,
\[
  \Gamma\left(
    V,\varprojlim_k\RR^n f_*(\sh{F}_k)
  \right)
  =
  \varprojlim_k
  \Gamma\bigl(V,\RR^n f_*(\sh{F}_k)\bigr)
\]
por \sref[0]{0.3.2.6}.
De \sref{III.4.1.7} y \sref{III.1.4.11} se sigue que
$\vphi_n$ es un isomorfismo topológico.
Del mismo modo, de nuevo por \sref{III.1.4.11}, \sref{III.4.1.7}
muestra que el homomorfismo $\psi_{n,V}$ de \eqref{III.4.1.3.5} es un
isomorfismo.
Por la definición de
$\RR^n\widehat{f}_*(\widehat{\sh{F}})$, se sigue que $\psi_n$ es un
isomorfismo.
\end{env}

\begin{corollary}[4.1.9]
\label{III.4.1.9}
Bajo las hipótesis de \sref{III.4.1.4}, para todo subconjunto abierto
afín $V$ de $Y$, el homomorfismo canónico
\[
  \HH^n\bigl(
    \widehat{X}\cap f^{-1}(V),
    \widehat{\sh{F}}
  \bigr)
  \longrightarrow
  \Gamma\bigl(
    \widehat{Y}\cap V,
    \RR^n\widehat{f}_*(\widehat{\sh{F}})
  \bigr)
\]
es biyectivo.
\end{corollary}

\begin{env}[4.1.10]
\label{III.4.1.10}
\emph{Observación.}
Sea $f:X\to Y$ un morfismo de tipo finito entre esquemas ordinarios
noetherianos, y sea $\sh{F}$ un $\sh{O}_X$-módulo coherente cuyo
soporte es propio sobre $Y$ \sref[II]{II.5.4.10}.
Sabemos entonces \sref{III.3.2.4} que $\RR^n f_*(\sh{F})$ es un
$\sh{O}_Y$-módulo coherente para todo $n\geq0$.
Además, siempre podemos suponer que
\[
  \sh{F}=u_*(\sh{G}),
  \qquad
  \sh{G}=u^*(\sh{F}),
\]
donde $\sh{G}$ es un $\sh{O}_Z$-módulo coherente, $Z$ es un
subesquema cerrado adecuado de $X$ cuyo espacio subyacente es
$\Supp(\sh{F})$, y $u:Z\to X$ es la inyección canónica
\sref[I]{I.9.3.5}.
Si
\[
  \sh{G}_k
  =
  \sh{G}\otimes_{\sh{O}_Y}
  (\sh{O}_Y/\sh{J}^{k+1}),
\]
entonces
\[
  \sh{G}_k=u^*(\sh{F}_k),
  \qquad
  \RR^n f_*(\sh{F}_k)
  =
  \RR^n(f\circ u)_*(\sh{G}_k)
\]
y
\[
  \RR^n f_*(\sh{F})
  =
  \RR^n(f\circ u)_*(\sh{G})
\]
por \sref{III.1.3.4}; y, teniendo en cuenta \sref[I]{I.10.9.5},
\[
  \RR^n\widehat{f}_*(\widehat{\sh{F}})
  =
  \RR^n\widehat{(f\circ u)}_*(\widehat{\sh{G}}).
\]
Podemos, por tanto, aplicar \sref{III.4.1.5} a $\sh{G}$ y al morfismo
propio $f\circ u$.
Se sigue que, bajo estas hipótesis, los resultados de
\sref{III.4.1.5} siguen siendo válidos para $\sh{F}$ y $f$.
\end{env}


\subsection{Casos particulares y variantes}
\label{subsection:III.4.2}

La forma más útil del teorema de comparación \sref{III.4.1.5} es la
siguiente.

\begin{proposition}[4.2.1]
\label{III.4.2.1}
Sea $Y$ un esquema localmente noetheriano, sea $f:X\to Y$ un
morfismo propio, y sea $\sh{F}$ un $\sh{O}_X$-módulo coherente.
Entonces, para todo $y\in Y$ y todo $p$,
\[
  \bigl(\RR^p f_*(\sh{F})\bigr)_y
\]

\oldpage[III]{130}

es un $\sh{O}_y$-módulo de tipo finito y está, por tanto, separado
para la topología $\mathfrak{m}_y$-preádica.
Además, existe un isomorfismo topológico canónico
\[
  \widehat{\bigl(\RR^p f_*(\sh{F})\bigr)_y}
  \isoto
  \varprojlim_k
  \HH^p\left(
    f^{-1}(y),
    \sh{F}\otimes_{\sh{O}_Y}
    (\sh{O}_y/\mathfrak{m}_y^k)
  \right).
  \tag{4.2.1.1}
  \label{III.4.2.1.1}
\]
Aquí el primer término es la completación de
$\bigl(\RR^p f_*(\sh{F})\bigr)_y$ para la topología
$\mathfrak{m}_y$-preádica.
En el segundo término, para todo $k\geq0$, $f^{-1}(y)$ se considera
como el espacio subyacente del esquema
\[
  X\times_Y\Spec(\sh{O}_y/\mathfrak{m}_y^k)
\]
\sref[I]{I.3.6.1}.
\end{proposition}

\begin{proof}
Como $\sh{O}_y$ es un anillo local noetheriano y
$\bigl(\RR^p f_*(\sh{F})\bigr)_y$ es un $\sh{O}_y$-módulo de tipo
finito \sref{III.3.2.1}, la topología $\mathfrak{m}_y$-preádica sobre
este módulo está separada \sref[0]{0.7.3.5}.

Cuando $Y$ es noetheriano e $y$ es cerrado, las demás afirmaciones se
siguen de \sref{III.4.1.7}, reemplazando $Y$ por un entorno afín de
$y$ y tomando $Y'=\{y\}$; aquí utilizamos (G,~II,~4.9.1).

En general, pongamos
\[
  Y_1=\Spec(\sh{O}_y),
  \qquad
  X_1=X\times_Y Y_1,
  \qquad
  \sh{F}_1
  =
  \sh{F}\otimes_{\sh{O}_Y}\sh{O}_{Y_1},
\]
y sea
\[
  f_1=f\times 1_{Y_1}:X_1\longrightarrow Y_1.
\]
El esquema $Y_1$ es noetheriano, el morfismo $f_1$ es propio
\sref[II]{II.5.4.2}[iii], y $\sh{F}_1$ es coherente
\sref[0]{0.5.3.11}.
Sea $y_1$ el único punto cerrado de $Y_1$.
La proposición ya demostrada se aplica a $f_1$, $\sh{F}_1$ y
$y_1$.
Tenemos
\[
  \sh{O}_{y_1}=\sh{O}_y,
  \qquad
  f_1^{-1}(y_1)=f^{-1}(y)
\]
\sref[I]{I.3.6.5}.
Además, los esquemas
\[
  X\times_Y\Spec(\sh{O}_y/\mathfrak{m}_y^k)
  \quad\text{y}\quad
  X_1\times_{Y_1}\Spec(\sh{O}_y/\mathfrak{m}_y^k)
\]
se identifican canónicamente \sref[I]{I.3.3.9}, al igual que
\[
  \sh{F}_1
  \otimes_{\sh{O}_{Y_1}}
  (\sh{O}_y/\mathfrak{m}_y^k)
  \quad\text{y}\quad
  \sh{F}
  \otimes_{\sh{O}_Y}
  (\sh{O}_y/\mathfrak{m}_y^k)
\]
\sref[I]{I.9.1.6}.
Solo queda observar que
\[
  \RR^p f_{1*}(\sh{F}_1)
  \simeq
  \RR^p f_*(\sh{F})
  \otimes_{\sh{O}_Y}
  \sh{O}_{Y_1}.
\]
Esto se sigue de \sref{III.1.4.15}, pues el morfismo local
$\Spec(\sh{O}_y)\to Y$ es plano
\sref[0]{0.6.7.1}, \sref[I]{I.2.4.2}.
\end{proof}

El corolario siguiente utiliza la terminología de la teoría de la
dimensión (capítulo~IV) y no se aplicará antes de dicho capítulo.

\begin{corollary}[4.2.2]
\label{III.4.2.2}
Sea $Y$ un esquema localmente noetheriano, sea $f:X\to Y$ un
morfismo propio, sea $y\in Y$, y sea $r$ la dimensión de
$f^{-1}(y)$.
Entonces, para todo $\sh{O}_X$-módulo coherente $\sh{F}$, los haces
$\RR^p f_*(\sh{F})$ se anulan en un entorno de $y$ para todo
$p>r$.
\end{corollary}

\begin{proof}
Para todo $k$, tenemos
\[
  \HH^p\left(
    f^{-1}(y),
    \sh{F}\otimes_{\sh{O}_Y}
    (\sh{O}_y/\mathfrak{m}_y^k)
  \right)
  =
  0
\]
por (G,~II,~4.15.2).
De \sref{III.4.2.1} se sigue que la completación separada de
$\bigl(\RR^p f_*(\sh{F})\bigr)_y$ para la topología
$\mathfrak{m}_y$-preádica es nula.
Como esta topología está separada, tenemos también
\[
  \bigl(\RR^p f_*(\sh{F})\bigr)_y=0.
\]
La conclusión se sigue porque $\RR^p f_*(\sh{F})$ es coherente
\sref[0]{0.5.2.2}.
\end{proof}

\begin{env}[4.2.3]
\label{III.4.2.3}
El resultado \sref{III.4.2.1} se utiliza especialmente cuando $p=0$.
Obtenemos el corolario siguiente.
\end{env}

\begin{corollary}[4.2.4]
\label{III.4.2.4}
Bajo las hipótesis de \sref{III.4.2.1}, existe un isomorfismo
topológico canónico
\[
  \widehat{\bigl(f_*(\sh{F})\bigr)_y}
  \isoto
  \varprojlim_k
  \Gamma\left(
    f^{-1}(y),
    \sh{F}_y/\mathfrak{m}_y^k\sh{F}_y
  \right).
\]
\end{corollary}


\subsection{Teorema de conexidad de Zariski}
\label{subsection:III.4.3}
\label{III.4.3}

Los resultados de esta sección y de la siguiente generalizan teoremas
bien conocidos de Zariski, y todos pueden deducirse de
\sref{III.4.2.4}. Son consecuencias del teorema siguiente.

\begin{theorem}[4.3.1]
\label{III.4.3.1}
\textnormal{(Teorema de conexidad).}
Sea $Y$ un esquema localmente noetheriano y sea $f:X\to Y$ un
morfismo propio. Entonces
\[
  \sh{A}(X)=f_*(\sh{O}_X)
\]
es una $\sh{O}_Y$-álgebra coherente.

\oldpage[III]{131}

Sea $Y'$ el $Y$-esquema finito tal que
$\sh{A}(Y')=\sh{A}(X)$; está determinado salvo $Y$-isomorfismo
\sref[II]{II.1.3.1}, \sref[II]{II.6.1.3}.
Si $f'=\sh{A}(e)$ es el $Y$-morfismo $X\to Y'$ deducido del
isomorfismo identidad
\[
  e:\sh{A}(Y')\isoto\sh{A}(X)
\]
\sref[II]{II.1.2.7}, entonces $f'$ es propio,
$f'_*(\sh{O}_X)$ es isomorfo a $\sh{O}_{Y'}$, y toda fibra
$f'^{-1}(y')$, $y'\in Y'$, es conexa y no vacía.
\end{theorem}

\begin{proof}
Sea $g:Y'\to Y$ el morfismo estructural. Para demostrar que el
homomorfismo
\[
  \theta:\sh{O}_{Y'}\longrightarrow f'_*(\sh{O}_X)
\]
que interviene en la definición de $f'$ es biyectivo, basta, como
$Y'$ es afín sobre $Y$, demostrar que
\[
  g_*(\theta):
  g_*(\sh{O}_{Y'})
  \longrightarrow
  g_*\bigl(f'_*(\sh{O}_X)\bigr)
  =
  f_*(\sh{O}_X)
\]
es la identidad \sref[II]{II.1.4.2}. Esto se sigue de las
definiciones, pues
$g_*(\sh{O}_{Y'})=\sh{A}(Y')$ y
$f_*(\sh{O}_X)=\sh{A}(X)$.
La coherencia de $\sh{A}(X)$ es un caso particular del teorema de
finitud \sref{III.3.2.1}. Como $f$ es propio y $g$ es separado,
$f'$ es propio \sref[II]{II.5.4.3}[i]. para terminar la demostración de
\sref{III.4.3.1}, basta, pues, demostrar el
\end{proof}

\begin{corollary}[4.3.2]
\label{III.4.3.2}
Bajo las hipótesis de \sref{III.4.3.1}, supongamos además que
$f_*(\sh{O}_X)$ es isomorfo a $\sh{O}_Y$. Entonces toda fibra
$f^{-1}(y)$ de $f$ es conexa y no vacía.
\end{corollary}

\begin{proof}
La hipótesis de que $f_*(\sh{O}_X)$ es isomorfo a $\sh{O}_Y$ implica
ya que $f$ es dominante, luego sobreyectivo porque $f$ es una
aplicación cerrada. Como en \sref{III.4.2.1}, podemos reducirnos al
caso en que $y$ es cerrado en $Y$. Como $f^{-1}(y)$ es un espacio
noetheriano, tiene un número finito de componentes conexas; estas son
también las componentes conexas del espacio subyacente de la
completación $\widehat{X}$ a lo largo de $f^{-1}(y)$.

Si $Z_i$, $1\leq i\leq n$, son estas componentes conexas, entonces
\[
  \Gamma(\widehat{X},\sh{O}_{\widehat{X}})
  =
  \prod_{i=1}^{n}\Gamma(Z_i,\sh{O}_{\widehat{X}}),
\]
y ninguno de los anillos del miembro derecho es nulo, pues la sección
unidad no es nula en ningún punto de $\widehat{X}$. Aplicando
\sref{III.4.1.5} a $\sh{F}=\sh{O}_X$, cuya completación a lo largo de
$f^{-1}(y)$ es $\sh{O}_{\widehat{X}}$, se ve que
$\Gamma(\widehat{X},\sh{O}_{\widehat{X}})$ es isomorfo a la
completación $\mathfrak{m}_y$-ádica separada
$\widehat{\sh{O}}_y$ del anillo local $\sh{O}_y$. Es, por tanto, un
anillo local y no puede ser un producto directo de varios anillos no
nulos. Así, $n=1$.
\end{proof}

\begin{corollary}[4.3.3]
\label{III.4.3.3}
Bajo las hipótesis de \sref{III.4.3.1}, para todo $y\in Y$ el conjunto
de componentes conexas de la fibra $f^{-1}(y)$ está en correspondencia
biyectiva con el conjunto finito de puntos de la fibra
$g^{-1}(y)$, donde $g:Y'\to Y$ es el morfismo estructural; en otros
términos, corresponde al conjunto de ideales maximales de
$\bigl(f_*(\sh{O}_X)\bigr)_y$.
\end{corollary}

\begin{proof}
Como $Y'$ es finito sobre $Y$, el espacio $g^{-1}(y)$ es finito y
discreto \sref[II]{II.6.1.7}. Como
\[
  f^{-1}(y)=f'^{-1}\bigl(g^{-1}(y)\bigr),
\]
la afirmación se sigue de esta observación y de
\sref{III.4.3.1}.
\end{proof}

Esto da una interpretación notable del $Y$-esquema $Y'$ definido en
\sref{III.4.3.1}. La factorización
$f=g\circ f'$ del morfismo propio $f$ es análoga a la factorización
obtenida por K.~Stein para aplicaciones holomorfas de espacios
analíticos, y se llamará en adelante la \emph{factorización de Stein}
de $f$.

\begin{remark}[4.3.4]
\label{III.4.3.4}
Sea $k$ una extensión del cuerpo $\kres(y)$. Si
\[
  f^{-1}(y)\otimes_{\kres(y)}k
  =
  X\times_Y\Spec(k)
\]
es conexo, entonces $f^{-1}(y)$ también es conexo, pues es la imagen
del primero por un morfismo de proyección
\sref[I]{I.3.4.7}. Para un morfismo $f:X\to Y$ de esquemas y un punto
$y\in Y$, diremos que la fibra $f^{-1}(y)$ es
\emph{geométricamente conexa} si, para toda extensión $k$ de
$\kres(y)$, el esquema
$f^{-1}(y)\otimes_{\kres(y)}k=X\times_Y\Spec(k)$ es conexo.

\oldpage[III]{132}

Bajo las hipótesis de \sref{III.4.3.2}, la conclusión puede
reforzarse: las fibras $f^{-1}(y)$ son geométricamente conexas.
En efecto, para toda extensión $k$ de $\kres(y)$ existen un anillo
local noetheriano $A$ y un homomorfismo local
\[
  \varphi:\sh{O}_y\longrightarrow A
\]
tal que $A$ es un $\sh{O}_y$-módulo plano y el cuerpo residual de
$A$ es $\kres(y)$-isomorfo a $k$ \sref[0]{0.10.3.1}.
Sea $Y_1=\Spec(A)$, y sea $h:Y_1\to Y$ el morfismo local
correspondiente a $\varphi$, que lleva el único punto cerrado $y_1$ de
$Y_1$ a $y$ \sref[I]{I.2.4.1}. Pongamos
\[
  X_1=X\times_Y Y_1,
  \qquad
  f_1=f\times 1_{Y_1}.
\]
El esquema $Y_1$ es noetheriano, $f_1$ es propio
\sref[II]{II.5.4.2}, y $f_1^{-1}(y_1)$ es un
$\kres(y_1)$-esquema isomorfo a $X\times_Y\Spec(k)$. Basta, por
tanto, demostrar que
\[
  f_{1*}(\sh{O}_{X_1})=\sh{O}_{Y_1},
\]
de modo que \sref{III.4.3.2} pueda aplicarse a $f_1$. Ahora $g$ es
plano por \sref[I]{I.2.4.2}; por tanto, por \sref{III.1.4.15} con
$q=0$,
\[
  f_{1*}(\sh{O}_{X_1})
  =
  h^*\bigl(f_*(\sh{O}_X)\bigr)
  =
  h^*(\sh{O}_Y)
  =
  \sh{O}_{Y_1}.
\]

En el caso general de \sref{III.4.3.1}, el mismo argumento muestra,
con la notación de dicho teorema, que
\[
  f_{1*}(\sh{O}_{X_1})
  =
  h^*\bigl(g_*(\sh{O}_{Y'})\bigr),
\]
y que la factorización de Stein $f_1=g_1\circ f'_1$ de $f_1$
satisface
\[
  g_1=g\times 1_{Y_1}
\]
\sref[II]{II.1.5.2}, siendo el $Y_1$-esquema finito correspondiente
$Y'_1=Y'\times_Y Y_1$. Teniendo en cuenta la transitividad de las
fibras \sref[I]{I.3.6.4}, se sigue de
\sref{III.4.3.3} que el número de componentes conexas de
$f_1^{-1}(y_1)$ es el número de elementos de
\[
  g_1^{-1}(y_1)
  =
  g^{-1}(y)\otimes_{\kres(y)}k.
\]
Si $k$ se elige algebraicamente cerrado, este número es independiente
de la extensión algebraicamente cerrada elegida y es igual al
\emph{número geométrico de puntos} de $g^{-1}(y)$
\sref[I]{I.6.4.7}, o equivalentemente a
\[
  \sum_i[\kres(y'_i):\kres(y)]_s,
\]
donde los $y'_i$ recorren el conjunto finito $g^{-1}(y)$. Este número
se llama \emph{número geométrico de componentes conexas} de
$f^{-1}(y)$. Los cuerpos $\kres(y'_i)$ son precisamente los cuerpos
residuales del anillo semilocal $\bigl(f_*(\sh{O}_X)\bigr)_y$.
\end{remark}

\begin{proposition}[4.3.5]
\label{III.4.3.5}
Sean $X$ e $Y$ esquemas íntegros localmente noetherianos, y sea
$f:X\to Y$ un morfismo propio dominante. Para todo $y\in Y$, el
número de componentes conexas de $f^{-1}(y)$ es a lo sumo el número
de ideales maximales de la clausura entera $\sh{O}'_y$ de
$\sh{O}_y$ en el cuerpo de funciones racionales $\mathrm{R}(X)$.
\end{proposition}

\begin{proof}
Para todo abierto $U$ de $Y$,
\[
  \Gamma\bigl(U,f_*(\sh{O}_X)\bigr)
  =
  \Gamma\bigl(f^{-1}(U),\sh{O}_X\bigr)
\]
es la intersección de los anillos locales $\sh{O}_x$ para
$x\in f^{-1}(U)$ \sref[I]{I.8.2.1.1}. Se sigue inmediatamente que
$\bigl(f_*(\sh{O}_X)\bigr)_y$ es un subanillo de $\mathrm{R}(X)$ que
contiene a $\sh{O}_y$. Como $f_*(\sh{O}_X)$ es un
$\sh{O}_Y$-módulo coherente, $\bigl(f_*(\sh{O}_X)\bigr)_y$ es un
$\sh{O}_y$-módulo finito y, por tanto, está contenido en $\sh{O}'_y$.
Todo ideal maximal de tal anillo es la intersección de este anillo con
un ideal maximal de $\sh{O}'_y$
\cite[vol.~I, pp.~257 and 259]{I-13}, lo que demuestra la proposición.
\end{proof}

\begin{definition}[4.3.6]
\label{III.4.3.6}
Un anillo local íntegro se dice \emph{unirrama} si su clausura entera
es de nuevo un anillo local. Un punto $y$ de un esquema íntegro
$Y$ se dice \emph{unirrama} si el anillo local $\sh{O}_y$ es
unirrama; este es, en particular, el caso cuando $Y$ es normal en
$y$.
\end{definition}

Sea $A$ un anillo local íntegro y sea $K$ su cuerpo de fracciones. El
anillo $A$ es unirrama si y solo si todo subanillo $A_1$ de $K$ que
contenga a $A$ y sea una $A$-álgebra finita es local. En efecto, sea
$A'$ la clausura entera de $A$. Por el primer teorema de
Cohen--Seidenberg (Bourbaki, \emph{Alg\`ebre commutative}, capítulo~V,
\S2, n.º~1, teorema~1), todo ideal maximal de $A_1$ es la contracción
a $A_1$ de un ideal maximal de $A'$; así, si $A'$ es local, también
lo es $A_1$.

\oldpage[III]{133}

Recíprocamente, $A'$ es el límite inductivo de la familia dirigida de
$A$-subálgebras finitas $A_\alpha$ de $A'$. Si cada $A_\alpha$ es
local, entonces, para $A_\alpha\subset A_\beta$, el ideal maximal de
$A_\alpha$ es la contracción del ideal maximal de $A_\beta$, por el
mismo argumento; por tanto, $A'$ es local
\sref[0]{0.10.3.1.3}.

Observemos también que, si la completación de un anillo local
noetheriano $A$ es íntegra ---propiedad que se expresa diciendo que
$A$ es \emph{analíticamente íntegro}---, entonces $A$ es unirrama.
Sea $\mathfrak{m}$ el ideal maximal de $A$, sea $K$ su cuerpo de
fracciones, y sea $K'$ el cuerpo de fracciones de $\widehat{A}$;
entonces
\[
  K'=K\otimes_A\widehat{A}.
\]
Sea $A'_F$ una $A$-subálgebra finita de $K$. El subanillo $B_F$ de
$K'$ generado por $\widehat{A}$ y $A'_F$ es isomorfo a
\[
  A'_F\otimes_A\widehat{A};
\]
es un $\widehat{A}$-módulo finito y es la completación de $A'_F$
para la topología $\mathfrak{m}$-ádica
\sref[0]{0.7.3.3}, \sref[0]{0.7.3.6}. Como $A'_F$ es semilocal
(Bourbaki, \emph{Alg\`ebre commutative}, capítulo~IV, \S2, n.º~5,
corolario~3 de la proposición~9), mientras que su completación es
íntegra, $A'_F$ solo puede tener un ideal maximal $\mathfrak{m}'_F$, y
$\mathfrak{m}'_F\cap A=\mathfrak{m}$. Esto demuestra la afirmación.

\begin{corollary}[4.3.7]
\label{III.4.3.7}
Bajo las hipótesis de \sref{III.4.3.5}, supongamos que la clausura
algebraica de $\mathrm{R}(Y)$ en $\mathrm{R}(X)$ tiene grado separable
$n$ y que $y\in Y$ es unirrama. Entonces la fibra $f^{-1}(y)$ tiene a
lo sumo $n$ componentes conexas. En particular, si la clausura
algebraica de $\mathrm{R}(Y)$ en $\mathrm{R}(X)$ es radicial sobre
$\mathrm{R}(X)$, entonces $f^{-1}(y)$ es conexa.
\end{corollary}

\begin{proof}
Sea $\sh{O}'_y$ la clausura entera de $\sh{O}_y$, y sea
$\sh{O}''_y$ la clausura entera de $\sh{O}_y$ en $\mathrm{R}(X)$.
Esta última es también la clausura entera de $\sh{O}'_y$ en
$\mathrm{R}(X)$. Si $\sh{O}'_y$ es local, entonces $\sh{O}''_y$ es
un anillo semilocal que tiene a lo sumo $n$ ideales maximales
\cite[vol.~I, p.~289, Theorem~22]{I-13}.

Este corolario es esencialmente la forma en la que Zariski enuncia su
``teorema de conexidad'' para esquemas algebraicos.
\end{proof}

\begin{remark}[4.3.8]
\label{III.4.3.8}
Si, además de las hipótesis de \sref{III.4.3.7}, $Y$ es normal en
$y$, entonces la fibra $f^{-1}(y)$ es geométricamente conexa: con la
notación de \sref{III.4.3.4}, $g^{-1}(y)$ es un solo punto $y'$, y
$\kres(y')$ es radicial sobre $\kres(y)$.
\end{remark}

\begin{definition}[4.3.9]
\label{III.4.3.9}
Sea $Y$ un esquema localmente noetheriano. Un morfismo de tipo finito
$f:X\to Y$ se dice \emph{universalmente abierto} si, para todo esquema
irreducible localmente noetheriano $Y'$ y todo morfismo dominante
$g:Y'\to Y$, toda componente irreducible de
$X'=X\times_Y Y'$ domina a $Y'$.
\end{definition}

Si $Y$ es irreducible, esto significa que, si $\eta$ y $\eta'$ son
los puntos genéricos de $Y$ e $Y'$, respectivamente, y si
$f'=f\times 1_{Y'}$, toda componente irreducible de $X'$ tiene su
punto genérico en $f'^{-1}(\eta')$. En consecuencia, para todo
abierto $U$ de $Y$, el morfismo restringido $f^{-1}(U)\to U$ es
universalmente abierto.

\begin{corollary}[4.3.10]
\label{III.4.3.10}
Sean $X$ e $Y$ esquemas íntegros localmente noetherianos, y sea
$f:X\to Y$ propio, dominante y universalmente abierto. Si la clausura
algebraica de $\mathrm{R}(Y)$ en $\mathrm{R}(X)$ es radicial sobre
$\mathrm{R}(Y)$, entonces toda fibra $f^{-1}(y)$, $y\in Y$, es
geométricamente conexa.
\end{corollary}

\begin{proof}
Podemos reducirnos al caso $Y=\Spec(B)$, donde $B$ es un anillo
íntegro noetheriano. Por \sref[II]{II.7.1.7}, existe un anillo local
noetheriano íntegramente cerrado $A$ que domina a $\sh{O}_y$ y tiene a
$\mathrm{R}(Y)$ como cuerpo de fracciones. Pongamos $Y'=\Spec(A)$, y
sea $h:Y'\to Y$ el morfismo correspondiente a la inyección canónica
$B\to A$. Es birracional, luego dominante; además, si $y'$ es el
único punto cerrado de $Y'$, entonces $h(y')=y$.

Pongamos
\[
  X'=X\times_Y Y',
  \qquad
  f'=f\times 1_{Y'}.
\]
Sean $\eta$, $\eta'$ y $\xi$ los puntos genéricos de $Y$, $Y'$ y
$X$, respectivamente. Entonces $f(\xi)=\eta$ y
$h^{-1}(\eta)=\{\eta'\}$; además,
$\kres(\eta)=\kres(\eta')=\mathrm{R}(Y)$. Así, $f'^{-1}(\eta')$ es
isomorfo a $f^{-1}(\eta)$ \sref[I]{I.3.6.4}, y tiene un único punto
genérico, pues $\xi$ es el punto genérico de
$f^{-1}(\eta)$ \sref[0]{0.2.1.8}. Por la apertura universal, toda
componente irreducible de $X'$ tiene su punto genérico en
$f'^{-1}(\eta')$. Por tanto, $X'$ es irreducible; su punto genérico
$\xi'$ es el de $f'^{-1}(\eta')$, y
$\kres(\xi')=\kres(\xi)$.

\oldpage[III]{134}

Pongamos $X''=X'_{\mathrm{red}}$ y $f''=f'_{\mathrm{red}}$. Entonces
$X''$ es íntegro y noetheriano, $f''$ es propio
\sref[II]{II.5.4.6}, y los espacios subyacentes de las fibras
$f'^{-1}(y')$ y $f''^{-1}(y')$ son iguales. Además,
\[
  \mathrm{R}(X'')=\kres(\xi')=\mathrm{R}(X).
\]
Así, $f''$ satisface las hipótesis de \sref{III.4.3.8}, y
$f''^{-1}(y')$ es geométricamente conexa.

Sea ahora $k$ una extensión cualquiera de $\kres(y)$. Existe una
extensión $k_1$ de $\kres(y)$ en la que tanto $\kres(y')$ como $k$
pueden considerarse subextensiones (Bourbaki, \emph{Alg\`ebre},
capítulo~V, \S4, proposición~2). Por el párrafo anterior,
\[
  f''^{-1}(y')\times_{Y'}\Spec(k_1)
\]
es conexo. Tiene el mismo esquema reducido asociado que
$f'^{-1}(y')\times_{Y'}\Spec(k_1)$
\sref[I]{I.5.1.8}, luego este último es conexo; es isomorfo a
$f^{-1}(y)\times_Y\Spec(k_1)$ \sref[I]{I.3.6.4}. Por tanto,
$f^{-1}(y)\times_Y\Spec(k)$ es conexo, y la observación al comienzo
de \sref{III.4.3.4} termina la demostración.
\end{proof}

\begin{remark}[4.3.11]
\label{III.4.3.11}
\emph{(i)}
El argumento anterior se debe, en lo esencial, a Zariski
\cite{I-20}, salvo que, en su marco, puede tomarse como $A$ la
clausura entera de $\sh{O}_y$, pues esta es un anillo noetheriano para
los anillos locales de la geometría algebraica clásica. Zariski
demuestra también que, si $Y$ es la variedad de Chow de un espacio
proyectivo $\mathbf{P}_k^r$ sobre un cuerpo $k$, y $X$ es el
subconjunto cerrado de $\mathbf{P}_k^r\times_kY$ que define la
correspondencia de Chow entre $\mathbf{P}_k^r$ e $Y$, entonces la
proyección $X\to Y$ es universalmente abierta (\emph{loc.\ cit.},
lema de la p.~82).
Parece que esta es la única propiedad formal de las ``coordenadas de
Chow'' utilizada en ciertas aplicaciones. En tal situación es, por
tanto, útil reemplazar el lenguaje de la especialización de ciclos en
el espacio proyectivo por el de las fibras de un morfismo propio,
supuesto eventualmente universalmente abierto o sometido a condiciones
análogas de regularidad local.

\emph{(ii)}
En el capítulo~IV veremos que un morfismo universalmente abierto
$f:X\to Y$ puede definirse equivalentemente exigiendo que, para todo
morfismo $Z\to Y$, el morfismo obtenido por cambio de base
$f_{(Z)}:X_{(Z)}\to Z$ sea abierto. También puede demostrarse que, si
$f$ satisface las hipótesis de \sref{III.4.3.10}, y si $y$ es una
especialización de $y'$, entonces el número geométrico de componentes
conexas de $f^{-1}(y)$ es a lo sumo el de $f^{-1}(y')$.
\end{remark}

\begin{corollary}[4.3.12]
\label{III.4.3.12}
Bajo las hipótesis de \sref{III.4.3.5}, supongamos además que
$\mathrm{R}(Y)$ es algebraicamente cerrado en $\mathrm{R}(X)$, y sea
$y$ un punto normal de $Y$. Entonces $f^{-1}(y)$ es geométricamente
conexa, y existe un entorno abierto $U$ de $y$ en $Y$ tal que
\[
  f_*\bigl(\sh{O}_X|_{f^{-1}(U)}\bigr)
  \simeq
  \sh{O}_Y|_U.
\]
En particular, si $Y$ es normal y $\mathrm{R}(Y)$ es algebraicamente
cerrado en $\mathrm{R}(X)$, entonces
$f_*(\sh{O}_X)\simeq\sh{O}_Y$.
\end{corollary}

\begin{proof}
La afirmación relativa a $f^{-1}(y)$ es un caso particular de
\sref{III.4.3.8}.

\oldpage[III]{135}

Si
\[
  X\xrightarrow{f'}Y'\xrightarrow{g}Y
\]
es la factorización de Stein de $f$ \sref{III.4.3.3}, se sigue que
$g^{-1}(y)$ consta de un solo punto $y'$. Además,
\[
  \sh{O}_y
  \subset
  \sh{O}_{y'}
  =
  \bigl(f_*(\sh{O}_X)\bigr)_y
  \subset
  \mathrm{R}(X).
\]
Como $\sh{O}_{y'}$ es finito sobre $\sh{O}_y$ (y \emph{a fortiori} sobre $\mathrm{R}(Y)$), está contenido en $\mathrm{R}(Y)$ en virtud de la hipótesis. Como $y$ es normal, necesariamente
$\sh{O}_{y'}=\sh{O}_y$. Así, $g$ es un isomorfismo local en $y'$
\sref[I]{I.6.5.4}, lo que demuestra la primera parte. La segunda
afirmación se sigue de la primera: la hipótesis adicional hace que $g$
sea biyectivo y un isomorfismo local en un entorno de cada punto de
$Y'$, luego un isomorfismo.
\end{proof}

El hecho de que \sref{III.4.3.7} esté disponible para esquemas da, por
ejemplo, la aplicación siguiente.

\begin{proposition}[4.3.13]
\label{III.4.3.13}
Sea $A$ un anillo local noetheriano unirrama, sea $\mathfrak{a}$ un
ideal de definición de $A$, y sean
\[
  A_0=A/\mathfrak{a},
  \qquad
  S=\operatorname{gr}_{\mathfrak{a}}(A)
\]
el anillo graduado asociado a la filtración $\mathfrak{a}$-preádica.
Supongamos que $S$ es una $A_0$-álgebra graduada generada por $S_1$,
donde $S_1$ es un $A_0$-módulo finito. Entonces $\Proj(S)$ es un
$A_0$-esquema conexo.
\end{proposition}

\begin{proof}
Sea $\mathfrak{m}$ el ideal maximal de $A$. El esquema
$Y=\Spec(A)$ es íntegro, y el punto $y$ correspondiente a
$\mathfrak{m}$ es su único punto cerrado. Para cierto entero $p$,
\[
  \mathfrak{m}^p\subset\mathfrak{a}\subset\mathfrak{m},
\]
de modo que $V(\mathfrak{a})=\{\mathfrak{m}\}$. Pongamos
\[
  S'=\bigoplus_{n\geq0}\mathfrak{a}^n
  \qquad\text{y}\qquad
  X=\Proj(S').
\]
Este es el $Y$-esquema obtenido por la explosión del ideal
$\mathfrak{a}$. Es íntegro, y el morfismo estructural
$f:X\to Y$ es birracional \sref[II]{II.8.1.4} y proyectivo.
Por tanto, \sref{III.4.3.7} se aplica y muestra que $f^{-1}(y)$ es
conexo. Pero el espacio subyacente de $f^{-1}(y)$ es el de
\[
  \Proj(S'\otimes_A A_0)
\]
\sref[I]{I.3.6.1}, \sref[II]{II.2.8.10}; como
$S'\otimes_AA_0=S$ por definición, se sigue la proposición.
\end{proof}


\subsection{El ``teorema principal'' de Zariski}
\label{subsection:III.4.4}

\begin{proposition}[4.4.1]
\label{III.4.4.1}
Sea $Y$ un esquema localmente noetheriano y sea $f:X\to Y$ un
morfismo propio. Sea $X'$ el conjunto de puntos $x\in X$ que están
aislados en su fibra $f^{-1}(f(x))$. Entonces $X'$ es abierto en $X$.
Si
\[
  f=g\circ f'
\]
es la factorización de Stein de $f$ \sref{III.4.3.3}, la restricción
de $f'$ a $X'$ es un isomorfismo de $X'$ sobre el subesquema inducido
en un subconjunto abierto $U$ de $Y'$, y
$X'=f'^{-1}(U)$.
\end{proposition}

\begin{proof}
Como $g^{-1}(f(x))$ es finito y discreto
\sref{III.4.3.3}, \sref[II]{II.6.1.7}, el punto $x$ está aislado en
$f^{-1}(f(x))$ si y solo si está aislado en
$f'^{-1}(f'(x))$. Podemos, por tanto, reducirnos al caso $f'=f$, de
modo que $f_*(\sh{O}_X)=\sh{O}_Y$. Si $x\in X'$, la fibra
$f^{-1}(f(x))$, que es conexa por \sref{III.4.3.2}, debe constar
entonces del único punto $x$.

Como $f$ es cerrado, para todo entorno abierto $V$ de $x$ en $X$, el
conjunto cerrado $f(X-V)$ no contiene a $y=f(x)$. Si $U$ es su
complementario en $Y$, entonces $f^{-1}(U)\subset V$. Así, las
imágenes inversas por $f$ de un sistema fundamental de entornos
abiertos de $y$ forman un sistema fundamental de entornos abiertos de
$x$. La igualdad $f_*(\sh{O}_X)=\sh{O}_Y$ y la definición de la
imagen directa de un haz \sref[0]{0.3.4.1}, \sref{III.4.2.1} implican
ahora que, si $f=(\psi,\theta)$, el homomorfismo
\[
  \theta_y^\#:\sh{O}_y\longrightarrow\sh{O}_x
\]
es un isomorfismo. Por tanto, existen entornos abiertos $V$ de $x$ y
$U$ de $y$ tales que la restricción de $f$ es un isomorfismo
$V\isoto U$ \sref[I]{I.6.5.4}.

\oldpage[III]{136}

Por el párrafo anterior podemos disponer además que
$f^{-1}(U)=V$. Se sigue inmediatamente que $V\subset X'$, lo que
termina la demostración.
\end{proof}

La proposición siguiente fue demostrada por Chevalley para esquemas
algebraicos.

\begin{proposition}[4.4.2]
\label{III.4.4.2}
Sea $Y$ un esquema localmente noetheriano y sea $f:X\to Y$ un
morfismo. Las condiciones siguientes son equivalentes:
\begin{enumerate}
  \item[{\rm a)}] $f$ es finito;
  \item[{\rm b)}] $f$ es afín y propio;
  \item[{\rm c)}] $f$ es propio y, para todo $y\in Y$, el conjunto
    $f^{-1}(y)$ es finito.
\end{enumerate}
\end{proposition}

\begin{proof}
La condición~a) implica b) por \sref[II]{II.6.1.2} y
\sref[II]{II.6.1.11}. Si $f$ es propio y afín, también lo es
\[
  f^{-1}(y)\longrightarrow\Spec(\kres(y))
\]
\sref[II]{II.1.6.2}, \sref[II]{II.5.4.2}. El teorema de finitud
\sref{III.3.2.1}, aplicado al haz estructural de $f^{-1}(y)$, muestra
que $f^{-1}(y)=\Spec(A)$ para una $\kres(y)$-álgebra finita $A$. Así,
$f^{-1}(y)$ es un conjunto finito \sref[II]{II.6.1.7}, y b) implica
c).

Por último, si se cumple c), la fibra $f^{-1}(y)$ es un esquema
algebraico sobre $\kres(y)$ cuyo conjunto subyacente es finito y es,
por tanto, discreto \sref[I]{I.6.4.4}. Con la notación de
\sref{III.4.4.1}, tenemos $X'=X$, luego
$f':X\to Y'$ es un isomorfismo. Como $g$ es finito, c) implica a).
\end{proof}

\begin{theorem}[4.4.3]
\label{III.4.4.3}
\textnormal{(``Teorema principal'' de Zariski).}
Sea $Y$ un esquema noetheriano, sea $f:X\to Y$ un morfismo
cuasiproyectivo, y sea $X'$ el conjunto de puntos $x\in X$ que están
aislados en su fibra $f^{-1}(f(x))$. Entonces $X'$ es abierto en $X$,
y el subesquema inducido por $X$ sobre $X'$ es isomorfo a un
subesquema inducido sobre un subconjunto abierto de un $Y$-esquema
$Y'$ finito sobre $Y$.
\end{theorem}

\begin{proof}
Existe un $Y$-esquema proyectivo $Z$ tal que $X$ es
$Y$-isomorfo al subesquema inducido sobre un subconjunto abierto de
$Z$ \sref[II]{II.5.3.2}, \sref[II]{II.5.5.1}. Nos reducimos así al
caso en que $f$ es proyectivo, luego propio \sref[II]{II.5.5.3}, y el
teorema se sigue inmediatamente de
\sref{III.4.4.1}.
\end{proof}

\begin{remark}[4.4.4]
\label{III.4.4.4}
Si $X$ es reducido (respectivamente, si $X$ es irreducible y $X'$ no
es vacío), puede exigirse en \sref{III.4.4.3} que $Y'$ sea reducido
(respectivamente, irreducible). En efecto, $Y'$ puede reemplazarse por
la clausura esquemática $\overline{X'}$ de $X'$ en $Y'$
\sref[I]{I.9.5.11}, \sref[II]{II.6.1.5}. Si $X'$ es reducido, también
lo es $\overline{X'}$ \sref[I]{I.9.5.9}; si $X'$ no es vacío y es
irreducible, entonces $\overline{X'}$ también es irreducible.
\end{remark}

\begin{corollary}[4.4.5]
\label{III.4.4.5}
Sea $Y$ un esquema localmente noetheriano, sea $f:X\to Y$ un
morfismo de tipo finito, y sea $x$ un punto de $X$ aislado en su fibra
$f^{-1}(f(x))$. Entonces existe un entorno abierto de $x$ en $X$
isomorfo a un subconjunto abierto de un $Y$-esquema finito sobre $Y$.
\end{corollary}

\begin{proof}
Pongamos $y=f(x)$, sea $U$ un entorno abierto afín de $y$ en $Y$, y
sea $V$ un entorno abierto afín de $x$ en $X$, contenido en
$f^{-1}(U)$. Como $Y$ es separado, la inclusión $U\to Y$ es afín
\sref[II]{II.1.6.3}; y como $V$ es afín sobre $U$ (\emph{ibid.}), la
restricción de $f$ a $V$ es un morfismo afín $V\to Y$
\sref[II]{II.1.6.2}. A fortiori, esta restricción es
cuasiproyectiva, pues es de tipo finito
\sref[I]{I.6.3.5}, \sref[II]{II.5.3.4}. Solo queda aplicar
\sref{III.4.4.3} a esta restricción.
\end{proof}

El Corolario \sref{III.4.4.5} puede enunciarse en el lenguaje del
álgebra conmutativa como sigue.

\begin{corollary}[4.4.6]
\label{III.4.4.6}
Sea $A$ un anillo noetheriano, sea $B$ una $A$-álgebra de tipo
finito, sea $\mathfrak q$ un ideal primo de $B$, y sea $\mathfrak p$
su imagen inversa en $A$. Supongamos que $\mathfrak q$ es a la vez
maximal y minimal en el conjunto de ideales primos de $B$ cuya imagen
inversa es $\mathfrak p$. Entonces existen un elemento
$g\in B-\mathfrak q$, una $A$-álgebra finita $A'$, y un elemento
$f'\in A'$ tales que las $A$-álgebras $B_g$ y $A'_{f'}$ son
isomorfas.
\end{corollary}

\begin{proof}
Apliquemos \sref{III.4.4.5} a $Y=\Spec(A)$ y $X=\Spec(B)$. La
hipótesis sobre $\mathfrak q$ dice precisamente que $\mathfrak q$
está aislado en su fibra \sref[I]{I.1.1.7}.
\end{proof}

La consecuencia siguiente parece menos general.

\oldpage[III]{137}

\begin{corollary}[4.4.7]
\label{III.4.4.7}
Sea $A$ un anillo local noetheriano, sea $B$ una $A$-álgebra de tipo
finito, y sea $\mathfrak n$ un ideal primo de $B$ cuya imagen inversa
en $A$ es el ideal maximal $\mathfrak m$. Supongamos que $\mathfrak n$
es maximal en $B$ y minimal en el conjunto de ideales primos de $B$
cuya imagen inversa es $\mathfrak m$ (lo que significa también que
$\mathfrak mB$ es $\mathfrak n$-primario). Entonces existen una
$A$-álgebra finita $A'$ y un ideal maximal $\mathfrak m'$ de $A'$
(cuya imagen inversa en $A$ es $\mathfrak m$) tales que $B_{\mathfrak
n}$ es isomorfo, como $A$-álgebra, a $A'_{\mathfrak m'}$.
\end{corollary}

El caso particular siguiente de \sref{III.4.4.7} se llama también a
veces el ``teorema principal''.

\begin{corollary}[4.4.8]
\label{III.4.4.8}
Bajo las hipótesis de \sref{III.4.4.7}, supongamos además que $A$ y
$B$ son dominios íntegros con el mismo cuerpo de fracciones $K$.
Si $A$ es íntegramente cerrado, entonces $A=B$.
\end{corollary}

\begin{proof}
La Observación \sref{III.4.4.4} muestra que, al aplicar
\sref{III.4.4.7}, podemos suponer que $A'$ es un dominio íntegro con
cuerpo de fracciones $K$. La hipótesis sobre $A$ da entonces
$A'=A$, luego $B_{\mathfrak n}=A$. Como
$A\subset B\subset B_{\mathfrak n}$, se sigue que $A=B$.
\end{proof}

El enunciado \sref{III.4.4.8} es la forma en la que Zariski dio su
``teorema principal'' (extendido aquí a anillos locales íntegros
noetherianos arbitrarios).

Los corolarios anteriores eran variantes locales de
\sref{III.4.4.3}, que es un resultado global. He aquí otra
consecuencia de naturaleza global.

\begin{corollary}[4.4.9]
\label{III.4.4.9}
Sea $Y$ un preesquema íntegro localmente noetheriano, y sea
$f:X\to Y$ un morfismo separado, de tipo finito y birracional.
Supongamos además que $Y$ es normal y que toda fibra $f^{-1}(y)$,
para $y\in Y$, es finita. Entonces $f$ es una inmersión abierta. Si,
además, $f$ es cerrado (en particular, si $f$ es propio), entonces
$f$ es un isomorfismo.
\end{corollary}

\begin{proof}
Sea $x\in X$ y pongamos $y=f(x)$. Como $f^{-1}(y)$ es un esquema
algebraico sobre $\kres(y)$, la hipótesis de que es finito implica que
es discreto \sref[I]{I.6.4.4}. Además, $\sh{O}_y$ es íntegramente
cerrado, y $\sh{O}_x$ y $\sh{O}_y$ tienen el mismo cuerpo de fracciones
\sref[I]{I.7.1.5}. Podemos, por tanto, aplicar
\sref{III.4.4.8}; si $f=(\psi,\theta)$, el homomorfismo
\[
  \theta_y^\#:\sh{O}_y\longrightarrow\sh{O}_x
\]
es biyectivo. De \sref[I]{I.6.5.4} se sigue que $f$ es un isomorfismo
local. Como $f$ es separado y $X$ es íntegro, $f$ es una inmersión
abierta \sref[I]{I.8.2.8}. La última afirmación se sigue del hecho de
que $f$ es dominante.
\end{proof}

\begin{proposition}[4.4.10]
\label{III.4.4.10}
Sea $Y$ un esquema localmente noetheriano, y sea $f:X\to Y$ un
morfismo localmente de tipo finito. El conjunto $X'$ de puntos
$x\in X$ que están aislados en su fibra $f^{-1}(f(x))$ es abierto en
$X$.
\end{proposition}

\begin{proof}
La cuestión es local sobre $X$ e $Y$, por lo que podemos suponer que
$X$ e $Y$ son afines noetherianos y que $f$ es de tipo finito. El
morfismo $f$ es entonces afín y de tipo finito, luego
cuasiproyectivo \sref[II]{II.5.3.4}; basta aplicar
\sref{III.4.4.3}.
\end{proof}

\oldpage[III]{138}

\begin{corollary}[4.4.11]
\label{III.4.4.11}
Sea $Y$ un esquema localmente noetheriano, y sea $f:X\to Y$ un
morfismo propio. El conjunto $U$ de puntos $y\in Y$ para los que
$f^{-1}(y)$ es discreto es abierto en $Y$, y el morfismo restringido
$f^{-1}(U)\to U$ es finito. En particular, todo morfismo propio
cuasifinito $X\to Y$ es finito.
\end{corollary}

\begin{proof}
El complementario de $U$ en $Y$ es la imagen por $f$ de $X-X'$, que
es cerrado en $X$ por \sref{III.4.4.10}; como $f$ es una aplicación
cerrada, $U$ es abierto. Además, \sref[II]{II.6.2.2} muestra que
$f^{-1}(y)$ es finita para todo $y\in U$. Como el morfismo restringido
$f^{-1}(U)\to U$ es propio \sref[II]{II.5.4.1}, es finito por
\sref{III.4.4.2}.
\end{proof}

\begin{remark}[4.4.12]
\label{III.4.4.12}
\emph{(i)}
Como se anunció en \sref[II]{II.6.2.7}, demostraremos en el capítulo~V
que, si $Y$ es localmente noetheriano, todo morfismo separado
cuasifinito $f:X\to Y$ es cuasiafín, luego cuasiproyectivo. Se seguirá
que la conclusión del teorema principal \sref{III.4.4.3} sigue siendo
válida cuando solo se supone que $f$ es separado y de tipo finito.
En efecto, \sref{III.4.4.10} muestra que $X'$ es abierto en $X$; y
como $X$ es localmente noetheriano, la restricción de $f$ a $X'$ sigue
siendo de tipo finito \sref[I]{I.6.3.5}, luego cuasifinita por la
definición de $X'$, y claramente separada. Podemos, por tanto, aplicar
\sref{III.4.4.3} a esta restricción.

\emph{(ii)}
En el capítulo~IV daremos una demostración más elemental de
\sref{III.4.4.10}, utilizando la teoría de la dimensión.
\end{remark}


\subsection{Completaciones de módulos de homomorfismos}
\label{subsection:III.4.5}

\begin{proposition}[4.5.1]
\label{III.4.5.1}
Sea $A$ un anillo noetheriano, sea $\mathfrak J$ un ideal de $A$, sea
$X$ un $A$-esquema de tipo finito, y sean $\sh{F}$ y $\sh{G}$ dos
$\sh{O}_X$-módulos coherentes cuyos soportes tienen intersección
propia sobre $Y=\Spec(A)$ \sref[II]{II.5.4.10}. Entonces, para todo
entero $n\geq0$,
\[
  \Ext^n_{\sh{O}_X}(X;\sh{F},\sh{G})
\]
es un $A$-módulo finito, y su completación separada para la topología
$\mathfrak J$-preádica se identifica canónicamente, con la notación de
\sref{III.4.1.7}, con
\[
  \Ext^n_{\sh{O}_{\hat{X}}}
    (\hat{X};\hat{\sh{F}},\hat{\sh{G}}).
\]
\end{proposition}

\begin{proof}
Sabemos por (T, 4.2) que existe una sucesión espectral birregular
$\mathrm{E}(\sh{F},\sh{G})$ cuyo término límite es
\[
  \Ext_{\sh{O}_X}(X;\sh{F},\sh{G})
\]
y cuyos términos $\mathrm{E}_2$ son
\[
  \mathrm{E}_2^{pq}
  =\HH^p\bigl(X,\shExt^q_{\sh{O}_X}(\sh{F},\sh{G})\bigr).
\]
El $\sh{O}_X$-módulo
$\shExt^q_{\sh{O}_X}(\sh{F},\sh{G})$ es coherente
\sref[0]{0.12.3.3}; su soporte está contenido en la intersección de
los soportes de $\sh{F}$ y $\sh{G}$ (T, 4.2.2) y, en consecuencia, es
propio sobre $Y$ \sref[II]{II.5.4.10}. De \sref{III.3.2.4} se sigue
que los $\mathrm{E}_2^{pq}$ son $A$-módulos finitos, y después, de
\sref[0]{0.11.1.8}, que lo mismo ocurre con todo término
$\mathrm{E}_r^{pq}$ de la sucesión espectral y con su término límite.

Por otra parte, si $i:\hat{X}\to X$ es el morfismo canónico, entonces
$\hat{\sh{F}}$ y $\hat{\sh{G}}$ se identifican canónicamente con
$i^*(\sh{F})$ e $i^*(\sh{G})$, y $i$ es plano
\sref[I]{I.10.8.8}, \sref[I]{I.10.8.9}. Para todo $q\geq0$ existe,
por tanto, un $i$-morfismo canónico
\[
  u_q:\shExt^q_{\sh{O}_X}(\sh{F},\sh{G})
  \longrightarrow
  \shExt^q_{\sh{O}_{\hat{X}}}
    (\hat{\sh{F}},\hat{\sh{G}})
\]
\sref[0]{0.12.3.4}, y el $\sh{O}_{\hat{X}}$-homomorfismo
correspondiente
\[
  u_q^\#:
  i^*\bigl(\shExt^q_{\sh{O}_X}(\sh{F},\sh{G})\bigr)
  \longrightarrow
  \shExt^q_{\sh{O}_{\hat{X}}}
    (\hat{\sh{F}},\hat{\sh{G}})
\]
es un isomorfismo \sref[0]{0.12.3.5}. En otros términos,
\sref[I]{I.10.8.8} identifica
\[
  \shExt^q_{\sh{O}_{\hat{X}}}
    (\hat{\sh{F}},\hat{\sh{G}})
\]
canónicamente con la completación de
$\shExt^q_{\sh{O}_X}(\sh{F},\sh{G})$, con la notación de
\sref{III.4.1.7}.

El teorema de comparación \sref{III.4.1.10} muestra ahora, para todo
$p\geq0$, que
\[
  \HH^p\bigl(\hat{X},
    \shExt^q_{\sh{O}_{\hat{X}}}
      (\hat{\sh{F}},\hat{\sh{G}})\bigr)
\]
se identifica canónicamente con la completación separada, para la
topología $\mathfrak J$-preádica, de
\[
  \HH^p\bigl(X,\shExt^q_{\sh{O}_X}(\sh{F},\sh{G})\bigr).
\]

\oldpage[III]{139}

Denotemos por $\mathrm{E}(\hat{\sh{F}},\hat{\sh{G}})$ la sucesión
espectral birregular de (T, 4.2) asociada a $\hat{\sh{F}}$ y
$\hat{\sh{G}}$. Si $\hat{A}$ es la completación separada de $A$ para
la topología $\mathfrak J$-preádica, entonces, salvo isomorfismo
canónico,
\[
  \mathrm{E}_2^{pq}(\hat{\sh{F}},\hat{\sh{G}})
  =
  \mathrm{E}_2^{pq}(\sh{F},\sh{G})\otimes_A\hat{A}
\]
\sref[0]{0.7.3.3}.

El morfismo plano $i$ define un homomorfismo canónico de sucesiones
espectrales
\[
  \varphi:\mathrm{E}(\sh{F},\sh{G})
  \longrightarrow
  \mathrm{E}(\hat{\sh{F}},\hat{\sh{G}})
  =
  \mathrm{E}(i^*(\sh{F}),i^*(\sh{G})).
\]
Sobre los términos $\mathrm{E}_2$ y sobre el término límite,
respectivamente, este es el homomorfismo
\[
  \varphi_2^{pq}:
  \HH^p\bigl(X,\shExt^q_{\sh{O}_X}(\sh{F},\sh{G})\bigr)
  \longrightarrow
  \HH^p\bigl(\hat{X},
    \shExt^q_{\sh{O}_{\hat{X}}}
      (\hat{\sh{F}},\hat{\sh{G}})\bigr)
\]
y el homomorfismo
\[
  \varphi^n:
  \Ext^n_{\sh{O}_X}(X;\sh{F},\sh{G})
  \longrightarrow
  \Ext^n_{\sh{O}_{\hat{X}}}
    (\hat{X};\hat{\sh{F}},\hat{\sh{G}}),
\]
deducidos de $u_q$ y $u_0$, respectivamente, por funtorialidad
\sref[0]{0.12.3.4}.

Tras tensorizar con $\hat{A}$, los $\varphi_r^{pq}$ y
$\varphi^n$ dan homomorfismos de $\hat{A}$-módulos
\[
  \psi_r^{pq}:
  \mathrm{E}_r^{pq}(\sh{F},\sh{G})\otimes_A\hat{A}
  \longrightarrow
  \mathrm{E}_r^{pq}(\hat{\sh{F}},\hat{\sh{G}})
\]
y
\[
  \psi^n:
  \Ext^n_{\sh{O}_X}(X;\sh{F},\sh{G})\otimes_A\hat{A}
  \longrightarrow
  \Ext^n_{\sh{O}_{\hat{X}}}
    (\hat{X};\hat{\sh{F}},\hat{\sh{G}}).
\]
Como $\hat{A}$ es un $A$-módulo plano \sref[0]{0.7.3.3}, los
$\hat{A}$-módulos
\[
  \mathrm{E}_r^{pq}(\sh{F},\sh{G})\otimes_A\hat{A}
\]
forman una sucesión espectral birregular cuyo término límite es
\[
  \Ext^n_{\sh{O}_X}(X;\sh{F},\sh{G})\otimes_A\hat{A},
\]
y los $\psi_r^{pq}$ y $\psi^n$ forman un morfismo de sucesiones
espectrales. Como los $\psi_2^{pq}$ son isomorfismos, lo mismo ocurre
con los $\psi^n$ \sref[0]{0.11.1.5}.
\end{proof}

\begin{corollary}[4.5.2]
\label{III.4.5.2}
Bajo las hipótesis de \sref{III.4.5.1}, supongamos además que $A$ es
un anillo noetheriano $\mathfrak J$-ádico. Entonces, para todo entero
$n\geq0$,
\[
  \Ext^n_{\sh{O}_X}(X;\sh{F},\sh{G})
  \isoto
  \Ext^n_{\sh{O}_{\hat{X}}}
    (\hat{X};\hat{\sh{F}},\hat{\sh{G}})
\]
canónicamente.
\end{corollary}

\begin{proof}
El $A$-módulo finito
$\Ext^n_{\sh{O}_X}(X;\sh{F},\sh{G})$ está separado y completo para la
topología $\mathfrak J$-preádica \sref[0]{0.7.3.6}.
\end{proof}

El caso particular $n=0$ de \sref{III.4.5.1} puede enunciarse como
sigue.

\begin{corollary}[4.5.3]
\label{III.4.5.3}
Bajo las hipótesis de \sref{III.4.5.1}, para todo homomorfismo
$u:\sh{F}\to\sh{G}$, denotemos por
$\hat{u}:\hat{\sh{F}}\to\hat{\sh{G}}$ el homomorfismo completado
\sref[I]{I.10.8.4}. Existe un isomorfismo canónico
\[
\label{III.4.5.3.1}
  \left(\Hom_{\sh{O}_X}(\sh{F},\sh{G})\right)^\wedge
  \isoto
  \Hom_{\sh{O}_{\hat{X}}}(\hat{\sh{F}},\hat{\sh{G}}),
  \tag{4.5.3.1}
\]
donde el miembro izquierdo es la completación separada, para la
topología $\mathfrak J$-preádica, del $A$-módulo
$\Hom_{\sh{O}_X}(\sh{F},\sh{G})$. Este isomorfismo se obtiene pasando
a las completaciones separadas en el homomorfismo
$u\mapsto\hat{u}$.
\end{corollary}


\subsection{Relaciones entre morfismos formales y morfismos ordinarios}
\label{subsection:III.4.6}

\begin{proposition}[4.6.1]
\label{III.4.6.1}
Sean $Y$ un esquema localmente noetheriano, $f:X\to Y$ un
morfismo propio, $\sh{F}$ un $\sh{O}_X$-módulo coherente que es
plano sobre $Y$, e $y\in Y$. Supongamos que, para algún entero
$n$,
\[
  \HH^n\bigl(f^{-1}(y),
    \sh{F}\otimes_{\sh{O}_Y}\kres(y)\bigr)=0.
\]
Entonces existe un entorno $U$ de $y$ en $Y$ tal que
\[
  \RR^n f_*(\sh{F})|U=0,
\]
y, para todo entero $p\geq0$, el homomorfismo canónico
\[
  \bigl(\RR^{n-1}f_*(\sh{F})\bigr)_y
  \longrightarrow
  \HH^{n-1}\bigl(f^{-1}(y),
    \sh{F}\otimes_{\sh{O}_Y}
      (\sh{O}_y/\mathfrak m_y^{p+1})\bigr)
\]
\eqref{III.4.2.1.1} es sobreyectivo.
\end{proposition}

\begin{proof}
El $\sh{O}_Y$-módulo $\RR^n f_*(\sh{F})$ es coherente
\sref{III.3.2.1}; por tanto, la primera afirmación se seguirá una vez que probemos
\[
  \bigl(\RR^n f_*(\sh{F})\bigr)_y=0
\]
\sref[0]{0.5.2.2}. Por \sref{III.4.2.1}, basta probar que
\[
  \HH^n\bigl(f^{-1}(y),
    \sh{F}\otimes_{\sh{O}_Y}
      (\sh{O}_y/\mathfrak m_y^{p+1})\bigr)=0
\]
para todo $p$. Esto es cierto por hipótesis para $p=0$; procedemos por
inducción sobre $p$.

Pongamos
\[
  X_p=X\times_Y
    \Spec(\sh{O}_y/\mathfrak m_y^{p+1}).
\]
Entonces $X_{p-1}$ es un subesquema cerrado de $X_p$ con el mismo
espacio subyacente \sref[I]{I.3.6.1}. La hipótesis de inducción
\[
  \HH^n\bigl(X_{p-1},
    \sh{F}\otimes_{\sh{O}_Y}
      (\sh{O}_y/\mathfrak m_y^p)\bigr)=0
\]
da, por tanto,
\[
  \HH^n\bigl(X_p,
    \sh{F}\otimes_{\sh{O}_Y}
      (\sh{O}_y/\mathfrak m_y^p)\bigr)=0.
\]
Por otra parte, la sucesión exacta de $\sh{O}_{X_p}$-módulos
\[
  0\longrightarrow
  \mathfrak m_y^p\sh{F}/\mathfrak m_y^{p+1}\sh{F}
  \longrightarrow
  \sh{F}/\mathfrak m_y^{p+1}\sh{F}
  \longrightarrow
  \sh{F}/\mathfrak m_y^p\sh{F}
  \longrightarrow0
\]
da la sucesión exacta de cohomología
\[
  \HH^n\bigl(X_p,
    \mathfrak m_y^p\sh{F}/\mathfrak m_y^{p+1}\sh{F}\bigr)
  \longrightarrow
  \HH^n\bigl(X_p,\sh{F}/\mathfrak m_y^{p+1}\sh{F}\bigr)
  \longrightarrow
  \HH^n\bigl(X_p,\sh{F}/\mathfrak m_y^p\sh{F}\bigr).
\]
Basta, pues, demostrar que
\[
\label{III.4.6.1.1}
  \HH^n\bigl(X_p,
    \mathfrak m_y^p\sh{F}/\mathfrak m_y^{p+1}\sh{F}\bigr)=0.
  \tag{4.6.1.1}
\]
En efecto, el término central será entonces un submódulo del último término,
que es cero por la hipótesis de inducción.

\oldpage[III]{140}

La fibra
\[
  Z=f^{-1}(y)=X\times_Y\Spec(\kres(y))
\]
es un subesquema cerrado de $X_p$, y
$\mathfrak m_y^p\sh{F}/\mathfrak m_y^{p+1}\sh{F}$ es anulado por
$\mathfrak m_y$. Por tanto, puede considerarse como un
$\sh{O}_Z$-módulo, y
\[
  \HH^n\bigl(Z,
    \mathfrak m_y^p\sh{F}/\mathfrak m_y^{p+1}\sh{F}\bigr)
  =
  \HH^n\bigl(X_p,
    \mathfrak m_y^p\sh{F}/\mathfrak m_y^{p+1}\sh{F}\bigr).
\]
Probaremos ahora que el $\sh{O}_Z$-homomorfismo canónico
\[
\label{III.4.6.1.2}
  (\sh{F}/\mathfrak m_y\sh{F})
    \otimes_{\kres(y)}
      (\mathfrak m_y^p/\mathfrak m_y^{p+1})
  \longrightarrow
  \mathfrak m_y^p\sh{F}/\mathfrak m_y^{p+1}\sh{F}
  \tag{4.6.1.2}
\]
es biyectivo. Una vez establecido esto, el hecho de que
$\mathfrak m_y^p/\mathfrak m_y^{p+1}$ sea un
$\kres(y)$-módulo libre da
\[
\begin{split}
  \HH^n\bigl(Z,
    \mathfrak m_y^p\sh{F}/\mathfrak m_y^{p+1}\sh{F}\bigr)
  &=
  \HH^n\bigl(Z,\sh{F}/\mathfrak m_y\sh{F}\bigr)
    \otimes_{\kres(y)}
      (\mathfrak m_y^p/\mathfrak m_y^{p+1})\\
  &=0
\end{split}
\]
\sref[0]{0.12.2.3}, puesto que
$\HH^n(Z,\sh{F}/\mathfrak m_y\sh{F})=0$ por hipótesis. Esto prueba
\eqref{III.4.6.1.1}.

Para terminar la demostración de la primera afirmación, solo queda demostrar
que \eqref{III.4.6.1.2} es biyectivo. La cuestión es puntual sobre
$X$, y $\sh{F}_x$ es un $\sh{O}_y$-módulo plano por hipótesis para
todo $x\in f^{-1}(y)$. Basta, pues, aplicar
\sref[0]{0.10.2.1}, c), puesto que
$\sh{F}_x/\mathfrak m_y\sh{F}_x$ es plano sobre el cuerpo
$\kres(y)=\sh{O}_y/\mathfrak m_y$.

Para la segunda afirmación de \sref{III.4.6.1}, nos reducimos de inmediato, como
en \sref{III.4.2.1}, al caso en que $Y$ es afín e $y$ es
cerrado. El mismo argumento demuestra, para todo $k>0$, que
\[
\label{III.4.6.1.3}
  \HH^n\bigl(X_{k+1},
    \mathfrak m_y^k\sh{F}/
      \mathfrak m_y^{k+p+1}\sh{F}\bigr)=0.
  \tag{4.6.1.3}
\]

\oldpage[III]{141}

De \sref{III.4.2.1} se sigue que
\[
\label{III.4.6.1.4}
  \bigl(\RR^n f_*(\mathfrak m_y^k\sh{F})\bigr)_y=0.
  \tag{4.6.1.4}
\]
La sucesión exacta de cohomología da ahora la exactitud de
\[
  \bigl(\RR^{n-1}f_*(\sh{F})\bigr)_y
  \longrightarrow
  \bigl(\RR^{n-1}f_*
    (\sh{F}/\mathfrak m_y^p\sh{F})\bigr)_y
  \longrightarrow
  \bigl(\RR^n f_*(\mathfrak m_y^p\sh{F})\bigr)_y=0.
\]
Como $y$ es cerrado e $Y$ es afín, tenemos
\sref{III.1.4.11}
\[
  \RR^{n-1}f_*(\sh{F}/\mathfrak m_y^p\sh{F})
  =
  \bigl(\HH^{n-1}
    (X,\sh{F}/\mathfrak m_y^p\sh{F})\bigr)^{\sim}
  =
  \bigl(\HH^{n-1}
    (f^{-1}(y),\sh{F}/\mathfrak m_y^p\sh{F})\bigr)^{\sim}
\]
(G, II, 4.9.1). El módulo
\[
  \HH^{n-1}
    (f^{-1}(y),\sh{F}/\mathfrak m_y^p\sh{F})
\]
es un $\sh{O}_y/\mathfrak m_y^p$-módulo, y por tanto
\[
  \bigl(\RR^{n-1}f_*
    (\sh{F}/\mathfrak m_y^p\sh{F})\bigr)_y
  =
  \HH^{n-1}
    (f^{-1}(y),\sh{F}/\mathfrak m_y^p\sh{F}).
\]
Esto termina la demostración.
\end{proof}

\begin{corollary}[4.6.2]
\label{III.4.6.2}
Sean $Y$ un esquema localmente noetheriano, $f:X\to Y$ un
morfismo propio y plano, $\sh{F}$ y $\sh{G}$ dos
$\sh{O}_X$-módulos localmente libres, e $y$ un punto de $Y$. Pongamos
\[
  X_y=f^{-1}(y)=X\times_Y\Spec(\kres(y)),\qquad
  \sh{F}_y=\sh{F}\otimes_{\sh{O}_Y}\kres(y),\qquad
  \sh{G}_y=\sh{G}\otimes_{\sh{O}_Y}\kres(y),
\]
y supongamos que
\[
\label{III.4.6.2.1}
  \HH^1\bigl(
    X_y,\shHom_{\sh{O}_{X_y}}(\sh{F}_y,\sh{G}_y)
  \bigr)=0.
  \tag{4.6.2.1}
\]
Entonces, para todo homomorfismo $u_0:\sh{F}_y\to\sh{G}_y$, existen
un entorno abierto $U$ de $y$ en $Y$ y un homomorfismo
\[
  u:\sh{F}|_{f^{-1}(U)}\longrightarrow
    \sh{G}|_{f^{-1}(U)}
\]
tal que $u_0=u\otimes1$.
\end{corollary}

\begin{proof}
La hipótesis nos permite aplicar \sref{III.4.6.1}, con $n=1$
y $p=0$, al $\sh{O}_X$-módulo coherente
\[
  \sh{H}=\shHom_{\sh{O}_X}(\sh{F},\sh{G}).
\]
En efecto, $\sh{H}$ es localmente libre y, por tanto, con mayor razón, plano sobre
$Y$, mientras que el $\sh{O}_{X_y}$-módulo
$\sh{H}\otimes_{\sh{O}_Y}\kres(y)$ se identifica con
$\shHom_{\sh{O}_{X_y}}(\sh{F}_y,\sh{G}_y)$
\sref[0]{0.6.2.2}. Podemos suponer que $Y=\Spec(A)$ es afín.
Entonces \sref{III.1.4.11} da
\[
  \RR^0f_*(\sh{H})
  =
  \bigl(\Hom_{\sh{O}_X}(\sh{F},\sh{G})\bigr)^{\sim},
\]
y, en consecuencia,
\[
  \bigl(\RR^0f_*(\sh{H})\bigr)_y
  =
  \Hom_{\sh{O}_X}(\sh{F},\sh{G})
    \otimes_A\sh{O}_y.
\]
Por \sref{III.4.6.1}, el homomorfismo canónico
\[
  \Hom_{\sh{O}_X}(\sh{F},\sh{G})
    \otimes_A\sh{O}_y
  \longrightarrow
  \Hom_{\sh{O}_{X_y}}(\sh{F}_y,\sh{G}_y)
\]
es sobreyectivo. Esto prueba el corolario, puesto que todo elemento del
módulo de la izquierda puede escribirse como $u\otimes(1/s)$, donde
$s\in A\setminus\mathfrak m_y$.
\end{proof}

\begin{corollary}[4.6.3]
\label{III.4.6.3}
Bajo las hipótesis de \sref{III.4.6.2}, si $u_0$ es inyectivo
(respectivamente sobreyectivo, biyectivo), se puede elegir $u$
inyectivo (respectivamente sobreyectivo, biyectivo).
\end{corollary}

\begin{proof}
Podemos restringirnos al caso $U=Y$. Basta probar
que, si $u_0$ es inyectivo (respectivamente sobreyectivo), entonces
\[
  \Ker u_x=0
  \quad\text{(respectivamente }\Coker u_x=0\text{)}
\]
para todo $x\in f^{-1}(y)$. En efecto, $\Ker u$ y $\Coker u$ son
$\sh{O}_X$-módulos coherentes \sref[0]{0.5.3.4}; por tanto, existe un
entorno $V$ de $f^{-1}(y)$ en $X$ sobre el cual la restricción de
$\Ker u$ (respectivamente $\Coker u$) es cero \sref[0]{0.5.2.2}.
Como $f$ es cerrado, existe un entorno abierto
$U'\subset U$ de $y$ tal que $f^{-1}(U')\subset V$, lo que prueba
la afirmación después de sustituir $U$ por $U'$.

\oldpage[III]{142}

Por hipótesis,
\[
  u_x\otimes1:
  \sh{F}_x\otimes_{\sh{O}_y}\kres(y)
  \longrightarrow
  \sh{G}_x\otimes_{\sh{O}_y}\kres(y)
\]
es inyectivo (respectivamente sobreyectivo), $\sh{F}_x$ y $\sh{G}_x$
son $\sh{O}_x$-módulos libres finitos, y $\sh{O}_x$ es plano sobre
$\sh{O}_y$. Si $u_x\otimes1$ es inyectivo, la inyectividad de $u_x$
se sigue de \sref[0]{0.10.2.4}. Si $u_x\otimes1$ es sobreyectivo,
entonces, con mayor razón, el homomorfismo inducido
\[
  \sh{F}_x/\mathfrak m_x\sh{F}_x
  \longrightarrow
  \sh{G}_x/\mathfrak m_x\sh{G}_x
\]
es sobreyectivo. Como $\sh{G}_x$ es un $\sh{O}_x$-módulo finito y
$\sh{O}_x$ es un anillo local de ideal maximal $\mathfrak m_x$, la
conclusión se sigue del lema de Nakayama (Bourbaki,
\emph{Alg\`ebre}, cap.~VIII, \S\,6, n.º~3, cor.~4 de la prop.~6).
\end{proof}

\begin{corollary}[4.6.4]
\label{III.4.6.4}
Sean $Y$ un esquema localmente noetheriano, $f:X\to Y$ un
morfismo propio y plano, e $y$ un punto de $Y$, y pongamos
$X_y=X\times_Y\Spec(\kres(y))$. Sea $\sh{E}_0$ un
$\sh{O}_{X_y}$-módulo localmente libre tal que
\[
\label{III.4.6.4.1}
  \HH^1\bigl(
    X_y,\shHom_{\sh{O}_{X_y}}(\sh{E}_0,\sh{E}_0)
  \bigr)=0.
  \tag{4.6.4.1}
\]
Sean $\sh{F}$ y $\sh{G}$ dos $\sh{O}_X$-módulos localmente libres
tales que $\sh{F}_y$ y $\sh{G}_y$, con la notación de
\sref{III.4.6.2}, son isomorfos a $\sh{E}_0$. Entonces existe un
entorno abierto $U$ de $y$ tal que
$\sh{F}|_{f^{-1}(U)}$ y $\sh{G}|_{f^{-1}(U)}$ son isomorfos.
\end{corollary}

\begin{corollary}[4.6.5]
\label{III.4.6.5}
Bajo las hipótesis de \sref{III.4.6.4} sobre $f$, $X$ e $Y$,
supongamos que
\[
  \HH^1(X_y,\sh{O}_{X_y})=0.
\]
Si $\sh{F}$ y $\sh{G}$ son dos $\sh{O}_X$-módulos invertibles tales
que $\sh{F}_y$ y $\sh{G}_y$ son isomorfos, entonces existe un entorno
abierto $U$ de $y$ tal que $\sh{F}|_{f^{-1}(U)}$ y
$\sh{G}|_{f^{-1}(U)}$ son isomorfos.
\end{corollary}

\begin{proof}
Basta aplicar \sref{III.4.6.4} a los módulos
$\sh{F}^{-1}\otimes\sh{G}$ y $\sh{O}_X$.
\end{proof}

\begin{remarks}[4.6.6]
\label{III.4.6.6}
\begin{enumerate}
\item[(i)] Utilizando \sref{III.4.6.5}, estableceremos en el capítulo~V
  la clasificación de los haces invertibles sobre un fibrado proyectivo
  anunciada en \sref[II]{II.4.2.7}.
\item[(ii)] El resultado de \sref{III.4.6.1} aparecerá en \S\,7 como
  consecuencia de proposiciones más generales.
\end{enumerate}
\end{remarks}

\begin{proposition}[4.6.7]
\label{III.4.6.7}
Sea $Z$ un esquema localmente noetheriano, y sean $X$ e $Y$
dos $Z$-esquemas cuyos morfismos estructurales
\[
  g:X\longrightarrow Z,\qquad h:Y\longrightarrow Z
\]
son propios. Sea $f:X\to Y$ un $Z$-morfismo, sea $z$ un punto
de $Z$, y pongamos
\[
  f_z=f\times_Z1:
  X\times_Z\Spec(\kres(z))
  \longrightarrow
  Y\times_Z\Spec(\kres(z)).
\]
\begin{enumerate}
\item[(i)] Si $f_z$ es finito (respectivamente una inmersión cerrada),
  existe un entorno abierto $U$ de $z$ tal que el morfismo
  \[
    g^{-1}(U)\longrightarrow h^{-1}(U)
  \]
  inducido por $f$ es finito (respectivamente una inmersión cerrada).
\item[(ii)] Supongamos además que $g$ es plano. Si $f_z$ es un
  isomorfismo, existe un entorno abierto $U$ de $z$ tal que
  el morfismo $g^{-1}(U)\to h^{-1}(U)$ inducido por $f$ es un
  isomorfismo.
\end{enumerate}
\end{proposition}

\begin{proof}
En ambos casos basta probar que, para todo
$y\in h^{-1}(z)$, existe un entorno $V_y$ de $y$ tal que
la restricción
\[
  f^{-1}(V_y)\longrightarrow V_y
\]
de $f$ es finita (respectivamente una inmersión cerrada, un isomorfismo).
En efecto, si $V$ es la unión de los $V_y$, entonces
$f^{-1}(V)\to V$ es finito (respectivamente una inmersión cerrada, un
isomorfismo) \sref[II]{II.6.1.1}, \sref[I]{I.4.2.4}. Como $h$ es
cerrado, existe un entorno abierto $U$ de $z$ tal que
$h^{-1}(U)\subset V$, y la proposición se sigue.

\noindent\textit{(i)} Obsérvese primero que $f$ es propio
\sref[II]{II.5.4.3}. Si
$f_z$ es finito, el entorno requerido $V_y$, para cada
$y\in h^{-1}(z)$, existe por \sref{III.4.4.11}.

\oldpage[III]{143}

Para tratar el caso en que $f_z$ es una inmersión cerrada, podemos
suponer ya, por tanto, que $f$ es finito. Así,
$X=\Spec(\sh{B})$, donde $\sh{B}$ es una
$\sh{O}_Y$-álgebra coherente, y $f$ corresponde \sref[II]{II.1.2.7} al
homomorfismo canónico
\[
  u:\sh{O}_Y\longrightarrow\sh{B}.
\]
Si, para todo $y\in h^{-1}(z)$, el homomorfismo
$u_y:\sh{O}_y\to\sh{B}_y$ es sobreyectivo, entonces $u$ es sobreyectivo en
algún entorno $V_y$ de $y$, puesto que $\Coker u$ es coherente
\sref[0]{0.5.3.4}, \sref[0]{0.5.2.2}. Ahora el morfismo finito
$f_z$ corresponde a
\[
  v=u\otimes1:
  \sh{O}_Y\otimes_{\sh{O}_Z}\kres(z)
  \longrightarrow
  \sh{B}\otimes_{\sh{O}_Z}\kres(z),
\]
y la hipótesis de que $f_z$ es una inmersión cerrada implica que
\[
  u_y\otimes1:
  \sh{O}_y\otimes_{\sh{O}_z}\kres(z)
  \longrightarrow
  \sh{B}_y\otimes_{\sh{O}_z}\kres(z)
\]
es sobreyectivo. Como $\sh{B}_y$ es un $\sh{O}_y$-módulo finito y
$\sh{O}_y$ es un anillo local noetheriano, la conclusión se sigue, como
en \sref{III.4.6.3}, del lema de Nakayama.

\noindent\textit{(ii)} El mismo razonamiento muestra que ahora
basta probar
que $u_y$ es biyectivo, sabiendo que $u_y\otimes1$ es biyectivo.
La inyectividad requerida se sigue del lema siguiente; la sobreyectividad
se sigue del argumento precedente.
\end{proof}

\begin{lemma}[4.6.7.1]
\label{III.4.6.7.1}
Sean $A$ y $B$ anillos locales noetherianos, sea
$\rho:A\to B$ un homomorfismo local, y sea $u:N\to M$ un
homomorfismo de $B$-módulos. Supongamos que $M$ es plano sobre $A$,
que $N$ es un $B$-módulo finito, y que
\[
  u\otimes1:
  N\otimes_A k\longrightarrow M\otimes_A k,
\]
donde $k$ es el cuerpo residual de $A$, es inyectivo. Entonces $N$ es
plano sobre $A$ y $u$ es inyectivo.
\end{lemma}

\begin{proof}
Para probar la primera afirmación, debemos demostrar que, para todo par de
$A$-módulos finitos $P,Q$ y todo $A$-homomorfismo inyectivo
$v:P\to Q$, el homomorfismo
\[
  1_N\otimes v:N\otimes_A P\longrightarrow N\otimes_A Q
\]
es inyectivo. Consideremos el diagrama conmutativo
\[
  \xymatrix{
    N\otimes_A P
      \ar[r]^-{1_N\otimes v}
      \ar[d]_{u\otimes1_P}
      &
    N\otimes_A Q
      \ar[d]^{u\otimes1_Q}
    \\
    M\otimes_A P
      \ar[r]_{1_M\otimes v}
      &
    M\otimes_A Q .
  }
\]
Como $1_M\otimes v$ es inyectivo por la planitud de $M$, basta
probar que $u\otimes1_P$ es inyectivo.

Sea $\mathfrak m$ el ideal maximal de $A$. La
filtración $\mathfrak m$-ádica sobre el $A$-módulo $N\otimes_A P$ es
también su filtración $\mathfrak mB$-ádica como $B$-módulo. La
topología definida por esta filtración es separada: $B$ es
noetheriano, $\mathfrak mB$ está contenido en el radical de $B$, y
$N\otimes_A P$ es un $B$-módulo finito, puesto que $N$ es finito sobre $B$
y $P$ es finito sobre $A$ \sref[0]{0.7.3.5}. Basta, por tanto,
probar que
\[
  \operatorname{gr}(u\otimes1_P):
  \operatorname{gr}_\bullet(N\otimes_A P)
  \longrightarrow
  \operatorname{gr}_\bullet(M\otimes_A P)
\]
es inyectivo, donde los módulos graduados asociados se toman con
respecto a las filtraciones $\mathfrak m$-ádicas (Bourbaki,
\emph{Alg\`ebre commutative}, cap.~III, \S\,2, n.º~8, cor.~1 del
teor.~1).

Como $M$ es plano sobre $A$, los homomorfismos
\[
  M\otimes_A(\mathfrak m^nP)
  \longrightarrow
  \mathfrak m^n(M\otimes_A P)
\]
son biyectivos. En consecuencia, el homomorfismo canónico
\[
  \vphi_M:
  \operatorname{gr}_0(M)\otimes_A\operatorname{gr}_\bullet(P)
  \longrightarrow
  \operatorname{gr}_\bullet(M\otimes_A P)
\]
es biyectivo.

\oldpage[III]{144}

Tenemos el diagrama conmutativo
\[
  \xymatrix{
    \operatorname{gr}_0(N)\otimes_A\operatorname{gr}_\bullet(P)
      \ar[r]^-{\operatorname{gr}_0(u)\otimes1}
      \ar[d]_{\vphi_N}
      &
    \operatorname{gr}_0(M)\otimes_A\operatorname{gr}_\bullet(P)
      \ar[d]^{\vphi_M}
    \\
    \operatorname{gr}_\bullet(N\otimes_A P)
      \ar[r]_{\operatorname{gr}(u\otimes1)}
      &
    \operatorname{gr}_\bullet(M\otimes_A P).
  }
\]
Aquí $\vphi_M$ es biyectivo y $\vphi_N$ es sobreyectivo.
Además, $\operatorname{gr}_0(u)$ es inyectivo por hipótesis, y
\[
\begin{split}
  \operatorname{gr}_0(N)\otimes_A\operatorname{gr}_\bullet(P)
    &=
  \operatorname{gr}_0(N)\otimes_k\operatorname{gr}_\bullet(P),\\
  \operatorname{gr}_0(M)\otimes_A\operatorname{gr}_\bullet(P)
    &=
  \operatorname{gr}_0(M)\otimes_k\operatorname{gr}_\bullet(P).
\end{split}
\]
Así, $\operatorname{gr}_0(u)\otimes1$ también es inyectivo. Del
diagrama se sigue que $\operatorname{gr}(u\otimes1_P)$ es
inyectivo, lo que prueba la primera afirmación. La segunda se sigue del
mismo argumento tomando $P=A$.
\end{proof}

\begin{proposition}[4.6.8]
\label{III.4.6.8}
Sea $Z$ un esquema localmente noetheriano, y sean $X$ e $Y$
dos $Z$-esquemas cuyos morfismos estructurales
$g:X\to Z$ y $h:Y\to Z$ son propios. Sea $Z'$ un subconjunto cerrado
de $Z$, pongamos
\[
  X'=g^{-1}(Z'),\qquad Y'=h^{-1}(Z'),
\]
y denotemos por
\[
  \widehat X=X_{/X'},\qquad
  \widehat Y=Y_{/Y'},\qquad
  \widehat Z=Z_{/Z'}
\]
las completaciones formales a lo largo de estos subconjuntos cerrados. Sea $f:X\to Y$
un $Z$-morfismo y sea
$\widehat f:\widehat X\to\widehat Y$ su extensión a las completaciones
formales. Para que $\widehat f$ sea un isomorfismo
(respectivamente una inmersión cerrada), es necesario y suficiente
que exista un entorno abierto $U$ de $Z'$ tal que el
morfismo
\[
  g^{-1}(U)\longrightarrow h^{-1}(U)
\]
inducido por $f$ sea un isomorfismo (respectivamente una inmersión cerrada).
\end{proposition}

\begin{proof}
La suficiencia de la condición es inmediata
\sref[I]{I.10.14.7}. Para probar su necesidad, basta, por el
mismo razonamiento que en \sref{III.4.6.7}, demostrar que para todo
$y\in Y'$ existe un entorno abierto $V_y$ de $y$ tal que
$f^{-1}(V_y)\to V_y$ sea un isomorfismo (respectivamente una
inmersión cerrada). Quedamos así reducidos al caso
\[
  Y=Z=\Spec(A),
\]
donde $A$ es noetheriano.

Por hipótesis \sref[I]{I.10.9.1}, \sref[I]{I.10.14.2}, la fibra
$f^{-1}(y)$ consta de un único punto para todo $y\in Y'$. Como
$f$ es propio \sref[II]{II.5.4.3}, \sref{III.4.4.11} da un entorno
abierto $U$ de $y$ para el cual $f^{-1}(U)\to U$ es finito. Podemos,
por tanto, suponer que $f$ es finito, de modo que
\[
  X=\Spec(B),
\]
donde $B$ es una $A$-álgebra finita.

Si $Y'=V(\mathfrak J)$, entonces
\[
  \widehat Y=\Spf(\widehat A),
  \qquad
  \widehat X=\Spf(\widehat B),
\]
donde $\widehat A$ es la completación separada de $A$ para la
topología $\mathfrak J$-ádica, y $\widehat B$ es la completación
separada de $B$ para la topología $\mathfrak JB$-ádica, o,
equivalentemente, la completación separada del $A$-módulo $B$ para
la topología $\mathfrak J$-ádica. Además, $\widehat f$
corresponde a la extensión continua
\[
  \widehat{\vphi}:\widehat A\longrightarrow\widehat B
\]
del homomorfismo canónico de anillos $\vphi:A\to B$, y la
hipótesis es que $\widehat{\vphi}$ es sobreyectivo (respectivamente
biyectivo) \sref[I]{I.10.14.2}. Pero $\widehat{\vphi}$ es también la
extensión continua de $\vphi$ considerado como un homomorfismo de
$A$-módulos. De \sref[I]{I.10.8.14} se sigue que existe un
entorno abierto $U$ de $Y'$ tal que la restricción a $U$ del
homomorfismo de $\sh{O}_Y$-módulos
\[
  \widetilde{\vphi}:\widetilde A\longrightarrow\widetilde B
\]
es sobreyectiva (respectivamente biyectiva), lo que termina la demostración.
\end{proof}


\subsection{Un criterio de amplitud}
\label{subsection:III.4.7}

\begin{theorem}[4.7.1]
\label{III.4.7.1}
Sean $Y$ un esquema localmente noetheriano, $f:X\to Y$ un
morfismo propio, $\sh{L}$ un $\sh{O}_X$-módulo invertible,
e $y$ un punto de $Y$. Pongamos
\[
  X_y=X\times_Y\Spec(\kres(y))=f^{-1}(y),
\]
y sea $g:X_y\to X$ el morfismo canónico. Si
\[
  \sh{L}_y=g^*(\sh{L})
  =\sh{L}\otimes_{\sh{O}_Y}\kres(y)
\]
es amplio sobre $X_y$, entonces existe un entorno abierto $U$ de $y$
en $Y$ tal que $\sh{L}|_{f^{-1}(U)}$ es amplio para la restricción
de $f$ a $f^{-1}(U)$.
\end{theorem}

\begin{proof}
\oldpage[III]{145}

\noindent\textit{I)} Pongamos
\[
  Y'=\Spec(\sh{O}_y),\qquad
  X'=X\times_Y Y',\qquad
  \sh{L}'=\sh{L}\otimes_{\sh{O}_Y}\sh{O}_y.
\]
Probamos primero que $\sh{L}'$ es amplio para $f'=f_{(Y')}$. Tenemos
el diagrama conmutativo
\[
  \xymatrix{
    X\ar[d]_f
      &
    X'\ar[l]\ar[d]_{f'}
      &
    X_y\ar[l]\ar[d]_{f_y}
    \\
    Y
      &
    Y'\ar[l]
      &
    \Spec(\kres(y))\ar[l].
  }
\]
Como $f'$ es propio \sref[II]{II.5.4.2}, (iii), y $\sh{O}_y$ es
noetheriano, \sref[II]{II.4.6.6} muestra que podemos reducirnos al caso
\[
  Y=Y'=\Spec(\sh{O}_y),\qquad X=X',
\]
suponer que $\sh{L}_y$ es amplio para $f_y$, y probar que $\sh{L}$
es amplio para $f$. Aplicamos el criterio
\sref{III.2.6.1}, c), y probamos que para todo
$\sh{O}_X$-módulo coherente $\sh{F}$ existe un entero $N$ tal que
\[
  \HH^1(X,\sh{F}(n))=0
  \quad\text{para todo }n\geq N,
  \qquad
  \sh{F}(n)=\sh{F}\otimes\sh{L}^{\otimes n}.
\]

El punto $y$ es cerrado en $Y$ y corresponde al ideal maximal
$\mathfrak m$ de $\sh{O}_y$. Por tanto, $X_y$ es el
subesquema cerrado de $X$ definido por el ideal coherente
\[
  \sh{J}
  =
  f^*(\widetilde{\mathfrak m})\sh{O}_X
  =
  \mathfrak m\sh{O}_X
\]
\sref[I]{I.4.4.5}, y $g$ es su inyección canónica. Consideremos
la $\kres(y)$-álgebra graduada
\[
  S=\operatorname{gr}(\sh{O}_y)
   =\bigoplus_{j\geq0}
      \mathfrak m^j/\mathfrak m^{j+1}.
\]
Es de tipo finito porque $\sh{O}_y$ es noetheriano. La
$\sh{O}_X$-álgebra
\[
  \sh{S}=f^*(\widetilde S)
\]
es cuasicoherente y de tipo finito, y es anulada por
$\sh{J}$. Así, si $\sh{S}_y=g^*(\sh{S})$, entonces $\sh{S}_y$ es una
$\sh{O}_{X_y}$-álgebra cuasicoherente de tipo finito y
\[
  \sh{S}=g_*(\sh{S}_y).
\]

Pongamos
\[
  \sh{M}_j
  =
  \mathfrak m^j\sh{F}/\mathfrak m^{j+1}\sh{F},
  \qquad
  \sh{M}
  =
  \bigoplus_{j\geq0}\sh{M}_j
  =
  \operatorname{gr}(\sh{F}).
\]
Como $\sh{F}$ es coherente, $\sh{M}$ es un
$\sh{S}$-módulo cuasicoherente de tipo finito \sref[0]{0.10.1.1}; también es
anulado por $\sh{J}$. En consecuencia, si
\[
  \sh{M}'_j=g^*(\sh{M}_j),
  \qquad
  \sh{M}'=g^*(\sh{M})
    =\bigoplus_{j\geq0}\sh{M}'_j,
\]
entonces $\sh{M}'$ es un $\sh{S}_y$-módulo graduado cuasicoherente de
tipo finito y $\sh{M}=g_*(\sh{M}')$. Además, poniendo
\[
  \sh{M}'_j(n)
  =
  \sh{M}'_j\otimes\sh{L}_y^{\otimes n},
\]
tenemos $\sh{M}'_j(n)=g^*(\sh{M}_j(n))$.

El morfismo $f_y$ es propio \sref[II]{II.5.4.2}, (iii), y
$\sh{L}_y$ es amplio; por tanto, $f_y$ es proyectivo
\sref[II]{II.5.5.4}, \sref[II]{II.4.6.11}. Podemos, pues, aplicar
\sref{III.2.4.1}, (ii), a
$\Spec(\kres(y))$, $f_y$, $\sh{S}_y$, $\sh{L}_y$ y $\sh{M}'$.
Existe un entero $N$ tal que, para $n\geq N$,
\[
  \HH^q(X_y,\sh{M}'_j(n))=0
\]
para todo $q>0$ y todo $j$. Se sigue que
\[
  \HH^q(X,\sh{M}_j(n))=0
\]
para todo $q>0$ y todo $j$ (G, II, 4.9.1).

Pongamos ahora
\[
  \sh{F}(n)_j
  =
  \sh{F}(n)/\mathfrak m^{j+1}\sh{F}(n).
\]
Para $j\geq1$ tenemos
\[
  \sh{F}(n)_{j-1}
  =
  \sh{F}(n)_j/\sh{M}_j(n),
  \qquad
  \sh{F}(n)_0=\sh{M}_0(n).
\]
Así, $\HH^1(X,\sh{F}(n)_0)=0$, y la sucesión exacta de cohomología
da, para todo $j\geq1$,
\[
  \HH^1(X,\sh{F}(n)_j)
  =
  \HH^1(X,\sh{F}(n)_{j-1}).
\]
Por tanto, $\HH^1(X,\sh{F}(n)_j)=0$ para todo $j\geq0$. Por
\sref{III.4.2.1}, se sigue que
\[
  \HH^1(X,\sh{F}(n))=0,
\]
lo que prueba la afirmación.

\oldpage[III]{146}

\noindent\textit{II)} Volvamos a la notación del comienzo de
la demostración. Podemos suponer que $Y=\Spec(A)$ es afín. Como
$f'$ es de tipo finito y $\sh{L}'$ es amplio para $f'$, existe un
entero $m>0$ tal que $\sh{L}'^{\otimes m}$ es muy amplio para
$f'$ \sref[II]{II.4.6.11}. Sustituyendo $\sh{L}$ por
$\sh{L}^{\otimes m}$, basta considerar el caso en que
$\sh{L}'$ es muy amplio para $f'$ y probar que
$\sh{L}|_{f^{-1}(U)}$ es muy amplio para la restricción de $f$.

Como $f'$ es propio, existe, para un entero adecuado $r>0$, una
$Y'$-inmersión cerrada
\[
  j:X'\longrightarrow P=\bb{P}_{Y'}^r
\]
tal que $\sh{L}'\simeq j^*(\sh{O}_P(1))$
\sref[II]{II.5.5.4}, (ii). Esta inmersión cerrada corresponde
canónicamente a un $\sh{O}_{X'}$-homomorfismo sobreyectivo
\[
  u:\sh{O}_{X'}^{r+1}\longrightarrow\sh{L}'
\]
\sref[II]{II.4.2.3}, o equivalentemente \sref[0]{0.5.1.1} a
$r+1$ secciones $s'_i$ $(0\leq i\leq r)$ de $\sh{L}'$ sobre $X'$
que generan este $\sh{O}_{X'}$-módulo.

Por definición, estas son también secciones de $f'_*(\sh{L}')$ sobre
$Y'$, y
\[
  f'_*(\sh{L}')
  =
  f_*(\sh{L})\otimes_{\sh{O}_Y}\sh{O}_{Y'}
\]
\sref[0]{0.5.4.10}. Como $Y$ es afín y $\sh{O}_y$ es el
anillo local de $A$ en el ideal primo $\mathfrak j_y$, cada $s'_i$
tiene la forma
\[
  s'_i=s''_i/t_i,
\]
donde $s''_i$ es una sección de $f_*(\sh{L})$ sobre $Y$ y
$t_i\in A\setminus\mathfrak j_y$. Se sigue que existen un
entorno abierto $V$ de $y$ en $Y$ y secciones $s_i$ de
$f_*(\sh{L})|_V$ tales que $s'_i=s_i/1$ (el espacio $Y'$ está
contenido en $V$; véase \sref[I]{I.2.4.2}).

Las secciones $s_i$ son secciones de $\sh{L}$ sobre $f^{-1}(V)$ y
definen, por tanto, un homomorfismo
\[
  v:
  \bigl(\sh{O}_X|_{f^{-1}(V)}\bigr)^{r+1}
  \longrightarrow
  \sh{L}|_{f^{-1}(V)}.
\]
Por construcción, es sobreyectivo en todo punto de $f^{-1}(y)$.
Como $\Coker(v)$ es coherente \sref[0]{0.5.3.4}, su soporte es
cerrado \sref[0]{0.5.2.2}; por tanto, existe un entorno abierto
$W\subset f^{-1}(V)$ de $f^{-1}(y)$ sobre el cual $v$ es sobreyectivo.
Como $f$ es cerrado, después de restringir alrededor de $y$ podemos elegir un
entorno abierto $U\subset V$ con $f^{-1}(U)\subset W$.
La conclusión se sigue ahora de \sref[II]{II.4.2.3}.
\end{proof}


\subsection{Morfismos finitos de esquemas formales}
\label{subsection:III.4.8}

\begin{proposition}[4.8.1]
\label{III.4.8.1}
Sea $\mathfrak{Y}$ un esquema formal localmente noetheriano,
$\sh{K}$ un ideal de definición de $\mathfrak{Y}$, y
$f:\mathfrak{X}\to\mathfrak{Y}$ un morfismo de esquemas formales.
Las condiciones siguientes son equivalentes:
\begin{enumerate}
  \item[\rm{a)}] $\mathfrak{X}$ es localmente noetheriano, $f$ es un
    morfismo ádico \sref[I]{I.10.12.1}, y, si
    $\sh{J}=f^*(\sh{K})\sh{O}_{\mathfrak{X}}$, el morfismo
    \[
      f_0:
      (\mathfrak{X},\sh{O}_{\mathfrak{X}}/\sh{J})
      \longrightarrow
      (\mathfrak{Y},\sh{O}_{\mathfrak{Y}}/\sh{K})
    \]
    inducido por $f$ es finito.
  \item[\rm{b)}] $\mathfrak{X}$ es localmente noetheriano y es el
    límite inductivo de un sistema inductivo ádico de $(Y_n)$, $(X_n)$,
    tal que el morfismo $X_0\to Y_0$ es finito.
  \item[\rm{c)}] Todo punto de $\mathfrak{Y}$ posee un entorno abierto
    formal afín noetheriano $V$ tal que $f^{-1}(V)$ es
    un abierto formal afín y
    $\Gamma(f^{-1}(V),\sh{O}_{\mathfrak{X}})$ es un
    $\Gamma(V,\sh{O}_{\mathfrak{Y}})$-módulo de tipo finito.
\end{enumerate}
\end{proposition}

\begin{proof}
De \sref[I]{I.10.12.3} se sigue inmediatamente que \emph{a)} implica
\emph{b)}. Para ver que \emph{b)} implica \emph{c)}, podemos suponer
que
\[
  \mathfrak{Y}=\Spf(B),\qquad \mathfrak{X}=\Spf(A),
\]
donde $B$ es un anillo noetheriano ádico y $\mathfrak{R}$ es un ideal
de definición de $B$. Por hipótesis, $X_0$ es un esquema afín
cuyo anillo $A_0$ es un $B/\mathfrak{R}$-módulo de tipo finito
\sref[II]{II.6.1.3}. Por \sref[I]{I.5.1.9}, cada $X_n$ es un esquema
afín; si $A_n$ denota su anillo, la condición \emph{b)} implica que,
para $m\leq n$,
\[
  A_m\simeq A_n/\mathfrak{R}^{m+1}A_n.
\]
Se sigue que $\mathfrak{X}\simeq\Spf(A)$, donde
$A=\varprojlim A_n$, y la conclusión se sigue de
\sref[0]{0.7.2.9}.

Finalmente, para probar que \emph{c)} implica \emph{a)}, podemos restringirnos
de nuevo al caso
\[
  \mathfrak{Y}=\Spf(B),\qquad \mathfrak{X}=\Spf(A),
\]
donde $A$ es una $B$-álgebra finita. Como
$A/\mathfrak{R}A$ es entonces una
$B/\mathfrak{R}$-álgebra finita, \sref[I]{I.10.10.9} muestra que se
satisfacen las condiciones de \emph{a)}.
\end{proof}

\begin{definition}[4.8.2]
\label{III.4.8.2}
Cuando se satisfacen las condiciones equivalentes \emph{a)}, \emph{b)} y \emph{c)}
de \sref{III.4.8.1}, se dice que el morfismo $f$ es
\emph{finito}; también se dice que $\mathfrak{X}$ es un
\emph{$\mathfrak{Y}$-esquema formal finito}, o un esquema
formal \emph{finito sobre $\mathfrak{Y}$}.
\end{definition}

\begin{proposition}[4.8.3]
\label{III.4.8.3}
\begin{enumerate}
  \item[\rm{(i)}] Una inmersión cerrada de esquemas formales localmente
    noetherianos es un morfismo finito.
  \item[\rm{(ii)}] El compuesto de dos morfismos finitos de esquemas formales
    localmente noetherianos es un morfismo finito.
  \item[\rm{(iii)}] Sean $\mathfrak{X}$, $\mathfrak{Y}$ y
    $\mathfrak{S}$ esquemas formales localmente noetherianos,
    $f:\mathfrak{X}\to\mathfrak{S}$ un morfismo finito, y
    $g:\mathfrak{Y}\to\mathfrak{S}$ un morfismo. Entonces la proyección
    \[
      \mathfrak{X}\times_{\mathfrak{S}}\mathfrak{Y}
      \longrightarrow\mathfrak{Y}
    \]
    es finita.
  \item[\rm{(iv)}] Sea $\mathfrak{S}$ un esquema formal localmente
    noetheriano, y sean $\mathfrak{X}'$ y $\mathfrak{Y}'$
    esquemas formales localmente noetherianos tales que
    $\mathfrak{X}'\times_{\mathfrak{S}}\mathfrak{Y}'$ es localmente
    noetheriano. Si $\mathfrak{X}$ y $\mathfrak{Y}$ son
    $\mathfrak{S}$-esquemas formales localmente noetherianos y
    \[
      f:\mathfrak{X}\to\mathfrak{X}',\qquad
      g:\mathfrak{Y}\to\mathfrak{Y}'
    \]
    son $\mathfrak{S}$-morfismos finitos, entonces
    $f\times_{\mathfrak{S}}g$ es finito.
  \item[\rm{(v)}] Sean
    \[
      f:\mathfrak{X}\to\mathfrak{Y},\qquad
      g:\mathfrak{Y}\to\mathfrak{Z}
    \]
    morfismos de esquemas formales localmente noetherianos tales que
    $g$ es de tipo finito y separado. Si $g\circ f$ es finito,
    entonces $f$ es finito.
\end{enumerate}
\end{proposition}

\begin{proof}
La afirmación \emph{(i)} es inmediata. Utilizando el criterio \emph{a)} de
\sref{III.4.8.1}, las demás afirmaciones se reducen de inmediato a las
proposiciones correspondientes para morfismos de esquemas ordinarios
\sref[II]{II.6.1.5}. Dejamos los detalles al lector, siguiendo
el modelo de \sref[I]{I.10.13.5}.
\end{proof}

\begin{corollary}[4.8.4]
\label{III.4.8.4}
Bajo las hipótesis de \sref[I]{I.10.9.9}, si $f$ es un morfismo
finito, también lo es su extensión $\widehat{f}$ a las
completaciones.
\end{corollary}

\begin{corollary}[4.8.5]
\label{III.4.8.5}
Sea $\mathfrak{X}$ un esquema formal finito sobre
$\mathfrak{Y}$, y sea
$f:\mathfrak{X}\to\mathfrak{Y}$ el morfismo estructural. Para
todo subconjunto abierto $U\subset\mathfrak{Y}$, el esquema formal
$f^{-1}(U)$ es finito sobre $U$.
\end{corollary}

\begin{proposition}[4.8.6]
\label{III.4.8.6}
Si $f:\mathfrak{X}\to\mathfrak{Y}$ es un morfismo finito de esquemas
formales localmente noetherianos, entonces
$f_*(\sh{O}_{\mathfrak{X}})$ es una
$\sh{O}_{\mathfrak{Y}}$-álgebra coherente.
\end{proposition}

\begin{proof}
Podemos considerar $f$ como el límite inductivo de un sistema inductivo
$(f_n)$ de morfismos $f_n:X_n\to Y_n$. Demostremos que los
$f_n$ son finitos y que $f_*(\sh{O}_{\mathfrak{X}})$ es
isomorfo al límite proyectivo de los
$(f_n)_*(\sh{O}_{X_n})$; esto probará la afirmación
\sref[I]{I.10.10.5}. Basta restringirse al caso
\[
  \mathfrak{Y}=\Spf(B),\qquad \mathfrak{X}=\Spf(A),
\]
y observar que, si $\mathfrak{R}$ es un ideal de definición
de $B$ y $A$ es un $B$-módulo de tipo finito, entonces
$A/\mathfrak{R}^{n+1}A$ es un módulo de tipo finito sobre
$B/\mathfrak{R}^{n+1}B$, y $A$ es el límite proyectivo de los
$A/\mathfrak{R}^{n+1}A$.
\end{proof}

Recíprocamente:

\begin{proposition}[4.8.7]
\label{III.4.8.7}
Sea $\mathfrak{Y}$ un esquema formal localmente noetheriano y
$\sh{A}$ una $\sh{O}_{\mathfrak{Y}}$-álgebra coherente. Existe
un esquema formal $\mathfrak{X}$ finito sobre $\mathfrak{Y}$,
único salvo $\mathfrak{Y}$-isomorfismo único, tal que
\[
  f_*(\sh{O}_{\mathfrak{X}})=\sh{A},
\]
donde $f:\mathfrak{X}\to\mathfrak{Y}$ es el morfismo estructural.
\end{proposition}

\begin{proof}
\oldpage[III]{148}
Sea $\sh{K}$ un ideal de definición de $\mathfrak{Y}$, y pongamos
\[
  Y_n=(\mathfrak{Y},
       \sh{O}_{\mathfrak{Y}}/\sh{K}^{n+1}),
  \qquad
  \sh{A}_n=\sh{A}/\sh{K}^{n+1}\sh{A}.
\]
Claramente, $\sh{A}_n$ es una $\sh{O}_{Y_n}$-álgebra finita y, por tanto,
define un $Y_n$-esquema finito
\[
  X_n=\Spec(\sh{A}_n)
\]
\sref[II]{II.6.1.3}. Para $m\leq n$, el homomorfismo canónico
sobreyectivo
\[
  h_{mn}:\sh{A}_n\longrightarrow\sh{A}_m
\]
define un morfismo $u_{mn}:X_m\to X_n$ para el cual el diagrama
\[
  \xymatrix{
    X_n\ar[d]_{f_n} &
    X_m\ar[l]_{u_{mn}}\ar[d]^{f_m}\\
    Y_n &
    Y_m\ar[l]
  }
\]
es conmutativo, donde $f_n$ denota el morfismo estructural; además,
identifica $X_m$ con
$X_n\times_{Y_n}Y_m$, como se sigue de inmediato de
\sref[II]{II.1.4.6}. El esquema formal $\mathfrak{X}$ obtenido
como límite inductivo del sistema inductivo $(X_n)$ es, por tanto,
localmente noetheriano, y su morfismo estructural
\[
  f:\mathfrak{X}\longrightarrow\mathfrak{Y},
\]
el límite inductivo del sistema $(f_n)$, es finito
\sref{III.4.8.1} \sref[I]{I.10.12.3.1}. Además, la demostración de
\sref{III.4.8.6} muestra que
$f_*(\sh{O}_{\mathfrak{X}})$ es el límite proyectivo de las
$\sh{A}_n$, y es, por tanto, igual a $\sh{A}$
\sref[I]{I.10.10.6}. La afirmación de unicidad se sigue del
resultado más general siguiente.
\end{proof}

\begin{proposition}[4.8.8]
\label{III.4.8.8}
Sea $\mathfrak{Y}$ un esquema formal localmente noetheriano, y
sean $\mathfrak{X}$ y $\mathfrak{X}'$
$\mathfrak{Y}$-esquemas formales finitos sobre $\mathfrak{Y}$, con
morfismos estructurales
\[
  f:\mathfrak{X}\to\mathfrak{Y},\qquad
  f':\mathfrak{X}'\to\mathfrak{Y}.
\]
Existe una biyección canónica
\[
  \Hom_{\mathfrak{Y}}(\mathfrak{X},\mathfrak{X}')
  \xrightarrow{\ \sim\ }
  \Hom_{\sh{O}_{\mathfrak{Y}}}
  \bigl(f'_*(\sh{O}_{\mathfrak{X}'}),
        f_*(\sh{O}_{\mathfrak{X}})\bigr)\footnotemark.
\]
\footnotetext{La última expresión denota el conjunto de homomorfismos de
$\sh{O}_{\mathfrak{Y}}$-álgebras
$f'_*(\sh{O}_{\mathfrak{X}'})\to f_*(\sh{O}_{\mathfrak{X}})$.}
\end{proposition}

\begin{proof}
La aplicación $h\mapsto\sh{A}(h)$ se define como en
\sref[II]{II.1.1.2}. Para probar que es biyectiva, nos reducimos de
inmediato al caso en que
$\mathfrak{Y}=\Spf(B)$ es un esquema formal afín noetheriano.
Entonces
\[
  \mathfrak{X}=\Spf(A),\qquad
  \mathfrak{X}'=\Spf(A'),
\]
donde $A$ y $A'$ son $B$-álgebras finitas, y
\[
  f_*(\sh{O}_{\mathfrak{X}})=A^\Delta,\qquad
  f'_*(\sh{O}_{\mathfrak{X}'})=A'^\Delta.
\]
La conclusión se sigue de la correspondencia biyectiva entre
$\mathfrak{Y}$-morfismos
$\mathfrak{X}\to\mathfrak{X}'$ y homomorfismos continuos de
$B$-álgebras $A'\to A$ \sref[I]{I.10.2.2}, junto con la
correspondencia biyectiva entre homomorfismos de $B$-módulos
$A'\to A$ y homomorfismos de
$\sh{O}_{\mathfrak{Y}}$-módulos
$A'^\Delta\to A^\Delta$ \sref[I]{I.10.10.2.3}.
\end{proof}

\begin{corollary}[4.8.9]
\label{III.4.8.9}
Bajo la biyección canónica de \sref{III.4.8.8}, las inmersiones
cerradas $\mathfrak{X}\to\mathfrak{X}'$ corresponden a
homomorfismos sobreyectivos de $\sh{O}_{\mathfrak{Y}}$-álgebras
\[
  f'_*(\sh{O}_{\mathfrak{X}'})
  \longrightarrow
  f_*(\sh{O}_{\mathfrak{X}}).
\]
\end{corollary}

\begin{proof}
La cuestión es local sobre $\mathfrak{Y}$ y, por tanto, se reduce a la
definición de las inmersiones cerradas de esquemas formales localmente
noetherianos \sref[I]{I.10.14.2}.
\end{proof}

\begin{corollary}[4.8.10]
\label{III.4.8.10}
Con la notación y las hipótesis de \sref{III.4.8.1}, un morfismo
ádico $f$ es una inmersión cerrada si y solo si $f_0$ es una
inmersión cerrada de esquemas ordinarios.
\end{corollary}

\begin{proof}
Esto se sigue de inmediato de \sref{III.4.8.9} y del criterio de
sobreyectividad para un homomorfismo de
$\sh{O}_{\mathfrak{Y}}$-módulos coherentes \sref[I]{I.10.11.5}.
\end{proof}

\begin{proposition}[4.8.11]
\label{III.4.8.11}
Un morfismo $f:\mathfrak{X}\to\mathfrak{Y}$ de esquemas formales
localmente noetherianos
\oldpage[III]{149}
es finito si y solo si es propio y todas sus fibras
$f^{-1}(y)$, para $y\in\mathfrak{Y}$, son finitas.
\end{proposition}

\begin{proof}
Por las definiciones \sref{III.3.4.1} y \sref{III.4.8.2}, la
afirmación se reduce inmediatamente a la misma afirmación para $f_0$
(con la notación de \sref{III.4.8.1}), que es precisamente
\sref{III.4.4.2}.
\end{proof}
